\documentclass[11pt]{article}
\usepackage{amsmath,amssymb,amsthm,mathtools,geometry,booktabs,enumitem}
\geometry{margin=1.1in}
\usepackage[colorlinks=true,linkcolor=blue,citecolor=blue]{hyperref}

\newtheorem{theorem}{Theorem}[section]
\newtheorem{lemma}[theorem]{Lemma}
\newtheorem{proposition}[theorem]{Proposition}
\newtheorem{corollary}[theorem]{Corollary}
\newtheorem{conjecture}[theorem]{Conjecture}
\theoremstyle{definition}
\newtheorem{definition}[theorem]{Definition}
\newtheorem{hypothesis}[theorem]{Hypothesis}
\newtheorem{openproblem}[theorem]{Open Problem}
\theoremstyle{remark}
\newtheorem{remark}[theorem]{Remark}

\newcommand{\A}{\mathcal{A}}
\newcommand{\w}{\omega_0}
\newcommand{\Hw}{\mathcal{H}_{\omega}}
\newcommand{\g}{\mathfrak{g}}
\newcommand{\E}{\mathcal{E}}
\newcommand{\Tr}{\operatorname{Tr}}
\newcommand{\W}{\mathbb{W}}

\title{The Physical State and Its Geometry:\\
\large a certified, gauge-free formulation of quantum matter,\\
from declared content to the continuum problem}
\author{Colin Bundschu}
\date{July 6, 2026}

\begin{document}
\maketitle

\begin{abstract}
A quantum theory of gauged matter is a single object, a state $\w$:
the minimizer of a linear energy functional over a unital $*$-algebra
of observables $\A$, an exact quotient with no gauge sector.  Every
quantity the theory emits is a certified two-sided bracket on a conic
linear program over the state's moment cone: a finite sum-of-squares
certificate below, an exhibited admissible state above, sound at every
stated budget independently of the convergence of any numerical
process.  The width of every emission is itself a read of the state,
its distance in the certificate's own seminorm from an exposed face of
the cone, and Heisenberg uncertainty and the exclusion principle are
the two graded instances of one positivity relation.

From the state we construct, canonically and without choices, its
carrier, closure ideal, modular structure, Dirichlet form and Dirac
operator, metric, and spectral dimension and volume: time, scale,
metric geometry, order, and frame are outputs of the state, not inputs
to the theory.  A field--gauge correspondence identifies compact-gauge
and CAR$\otimes$CCR content as two generating views of one dual-pair
carrier, the Gauss law their singlet compression; on it the group's
rank is an output of the closure ideal and, conditionally, the
signature of spacetime is a modular output.  Space is proposed as the
modular structure of a committed net built from the state's own
metric, with four decisive tests specified in advance and none yet
executed.

The thermodynamic limit is proved for energy densities, with a
finite-size gap criterion available on the stoquastic class, whose
boundary is the fermion sign problem, stated here as an open
structural question.  The continuum limit is reduced to a
certification problem: refinement acts on moment bodies by exact
linear projection, failure of certified contraction must localize at a
genuine zero-slack state, no accumulation of certified data can decide
between the branches, and, conditionally on a single parametric
certificate at the Gaussian fixed point and a finite chain of step
certificates, the continuum brackets of an asymptotically free theory
constitute a finite, machine-verifiable problem.  The existence of the
chain is the central conjecture: nonperturbative asymptotic freedom in
certificate form.

A fully worked instance, SU(2) gauge theory on two plaquettes,
supplies certified spectra, closed-form geometry, and coupling towers,
and passes the framework's first decisive test: its Weyl consistency
requirement, certified spectral volume against the Hausdorff volume of
the state's metric, closes exactly, both sides in closed form.
Electrodynamic matter, with gravitation carried as a spin-2 boson
sector, is developed as declared content of the same framework; its
reads recover spectroscopy, order, frame, response, and the measured
gravitational quantities, and an executed relational control in
hydrogen exercises the emission pipeline end to end on a case where
both numbers are independently known.
\end{abstract}

\section{Introduction}\label{sec:intro}

Three commitments define this formulation.

\emph{First, the state is the only unknown.}  A theory is declared by its
content: an algebra of observables presented as an exact quotient (the
gauge sector removed at the level of the algebra, never fixed), a
canonical family of derivations induced by the content's symmetry, and a
Hamiltonian element.  Everything else (the carrier Hilbert space, the
ground representative, the flow of time, the metric, the dimension, the
organization of matter) is constructed from the state or read from it.
Where the construction requires an unproven statement, that statement
appears below as a named hypothesis, conjecture, or open problem, in
place, with the results that depend on it stated conditionally.

\emph{Second, every emission is a certified two-sided bracket.}  The
instrument never emits an estimate.  A lower bound is a sum-of-squares
certificate over the algebra, verified as a finite exact identity; an
upper bound is the exact evaluation of an exhibited admissible state.
Both are valid at every stated budget regardless of how they were found;
no optimizer's convergence is load-bearing for truth.  Emitted quantities
are dimensionless and relational; dimensionful intermediates never leave
the instrument.

\emph{Third, nonlinearity in the state is the signature of a gauge
mode.}  Every physical functional in this paper is linear in the state,
an element of a convex cone, or an envelope of affine functions of the
state (the supporting-hyperplane structure of the moment cone made
visible); wherever a construction appeared to require a nonlinear object
(a wavefunction, a chart, a threshold, a rescaled block field), the
nonlinearity marked a representation artifact, and the construction is
given here in its linear form.

\subsection*{Vocabulary}
The paper uses a small fixed working vocabulary throughout; its nearest
standard equivalents are these.

\begin{center}
\begin{tabular}{@{}lp{0.66\textwidth}@{}}
\toprule
term & nearest standard reading \\
\midrule
content & a presented theory: algebra (exact quotient), canonical
derivations, Hamiltonian element, budget axes \\
budget & a declared finite truncation axis (word degree, spectral
cutoff, band, volume) \\
read & a derived quantity, linear (or an affine envelope) in the state \\
emission & a certified two-sided enclosure $[L,U]$ of a read \\
bracket & the interval $[L,U]$ itself \\
certificate & a finite sum-of-squares identity with multipliers,
verified in exact or interval arithmetic \\
width & $U - L$, itself a read of the state
(Theorem~\ref{thm:gapid}) \\
instrument & the certification pipeline as a whole \\
fails shut & returns no bound rather than an unsound one \\
rung & a level of the constructive trial ladder (\S\ref{sec:ladder}) \\
pin & an executed check against an independently known value \\
\bottomrule
\end{tabular}
\end{center}

\subsection*{Reading order}
Sections~\ref{sec:kinematics}--\ref{sec:canonical} are the framework
and are self-contained; a reader who wants the framework exercised can
pass from \S\ref{sec:kinematics} directly to the worked instance of
\S\ref{sec:instance} and its appendices.
\S\S\ref{sec:extensivity}--\ref{sec:continuum} carry the
complete-proof analysis of the thermodynamic and continuum questions;
\S\ref{sec:correspondence} is classical invariant theory and Howe
duality read into the frame; \S\ref{sec:space}, with the real-form
discussion it carries, is the conditional register of the paper,
committed constructions with named decisive tests and nothing
exhibited; \S\ref{sec:subsumption} is the framework evaluated at
electrodynamic content and can be read independently after
\S\ref{sec:kinematics}.  Proof statuses are graded in
\S\ref{sec:conclusion}.

The paper proves what can be proved and states precisely what cannot yet
be.  The main results are: the uncertainty hierarchy that identifies every
emitted width with the state's distance, in the certificate's own seminorm,
from an exposed face of the moment cone, and places Heisenberg uncertainty
and the exclusion principle as the two graded instances of one relation
(\S\ref{sec:uncertainty}); the constructive upper bracket, a single trial
ladder whose reference rung is the kinetic vacuum read at the content's
group: Haar for gauge content, quasi-free$\otimes$Gaussian for field
content, one construction at two groups (\S\ref{sec:ladder}); the
field--gauge correspondence, on which compact-gauge and field content
are two generating views of one dual-pair carrier, the Gauss law is
singlet compression, and the group's rank is an output of the state's
closure ideal, the spacetime group entering as its complex point with
the real form a modular output
(\S\S\ref{sec:correspondence},~\ref{sec:space}); the canonical
construction stack of \S\ref{sec:canonical}, including the automatic local
factorization of
stoquastic-class Hamiltonians in the ground carrier
(Theorem~\ref{thm:factorization}), with the fermion sign problem
identified as its exact scope boundary (Open Problem~\ref{op:sign}); the geometric identification and its
falsifiable Weyl consistency requirement, verified in closed form on
the worked instance (\S\ref{sec:geometry},
Theorem~\ref{thm:weylinstance}); complete
extensivity theorems for energy densities and a conditional finite-size
gap criterion (\S\ref{sec:extensivity}); the reduction of the continuum
problem to a finite certification problem, comprising the exactness of
refinement (Proposition~\ref{prop:projection}), the two-dimensional
theorem (Theorem~\ref{thm:2d}), the dichotomy and localization theorems
(Theorems~\ref{thm:dichotomy}, \ref{thm:totaldichotomy}), the no-go
theorem on data (Theorem~\ref{thm:nogo}), and the conditional Reduction
Theorem (Theorem~\ref{thm:reduction}), with the central conjecture
(Conjecture~\ref{conj:chain}) stated in minimal form
(\S\ref{sec:continuum}); the hadronic content and its certified targets
(\S\ref{sec:hadrons}); and electrodynamic matter developed as declared
content, with its single-emission structure, convergence theorems,
gravitational reads, and executed hydrogen control
(\S\ref{sec:subsumption}).  A worked
instance with certified numbers runs through the paper
(\S\ref{sec:instance}).

\section{Kinematics and emissions}\label{sec:kinematics}

\subsection{Content}\label{sec:content}
A theory is declared by: (i) a unital $*$-algebra $\A$, presented as an
exact quotient so that no gauge degree of freedom exists in $\A$; for
lattice gauge content with compact group $G$ on a declared face set, $\A$
is the algebra of invariants (generated by loop words), and the quotient
requires no gauge fixing and no section; (ii) the canonical derivations
$\{X_a\}$: the action of the Lie algebra $\g$ on $\A$ induced by the
content's group action, with kinetic form the Casimir of the action
(unique up to the normalization quotiented below); (iii) a coupling
$\lambda$ and a Hermitian magnetic element $M \in \A$, assembling the
Hamiltonian $h(\lambda) = \mathrm{Cas} - \lambda M$; (iv) budget axes:
the canonical filtrations of $\A$ by degree over the canonical generators
and by spectral cutoff of the Casimir.  Nothing else is input.  In
particular no Hilbert space, metric, dimension, lattice spacing,
coordinate, or basis is declared.  The group in (i) and (ii) enters
with the standing of a presentation, not of an independent input: on
the correspondence of \S\ref{sec:correspondence} the same content is
declared by its fusion functional alone, the pairing symmetry of the
words carrying the content and the rank recovered from the closure
ideal of the state (Proposition~\ref{prop:groupclosure}); the group is
a choice of generators, and emissions do not depend on it (exactly at
the worked content, Proposition~\ref{prop:prepot}; per Open
Problem~\ref{op:dualpair} in general).

\subsection{States, cones, certificates}\label{sec:cones}
States are normalized positive linear functionals on $\A$.  At budget $b$
a state restricts to its degree-$\le b$ moment data.  Two objects at
budget $b$ must be held apart.  The \emph{moment body} $\mathcal S_b$
is the set of such restrictions over admissible states.  The
\emph{moment cone} $\mathcal M_b$ is its computable outer
approximation: the convex
set of degree-$\le b$ moment data satisfying the linear matrix
inequalities expressing positivity on squares.  The inclusion
$\mathcal S_b \subseteq \mathcal M_b$ is generally strict at finite
$b$; its defect is exactly what Theorems~\ref{thm:gapid}
and~\ref{thm:closure} quantify, and conflating the two objects would
collapse the hierarchy.  Certificates are elements
of the dual cone: finite sum-of-squares identities with multipliers over
the declared constraint set.  An \emph{emission} is a certified interval
$[L,U] \ni v$ for a value
\[
v \;=\; \inf_{\omega \in \mathcal S} \,\omega(A)
\qquad (\text{or } \sup),
\]
a linear functional optimized over a convex set $\mathcal S$ of
admissible states: $L$ from a certificate by weak duality, $U$ from an
exhibited member of $\mathcal S$ evaluated exactly.

\begin{proposition}[Weak duality; anytime validity]\label{prop:weak}
If $A - \lambda\mathbf 1 - \sum_i \mu_i C_i = \sum_j S_j^\dagger S_j$
holds as a finite identity in $\A$, with $C_i$ the constraint elements
vanishing in expectation on $\mathcal S$, then $\omega(A) \ge \lambda$
for every $\omega \in \mathcal S$.  If the identity holds up to a
residual $R$ of bounded norm, then $\omega(A) \ge \lambda - \|R\|$.
Both statements hold for arbitrary coefficient tuples; consequently every
solver iterate emits a valid bound, and no convergence claim is
load-bearing.
\end{proposition}
\begin{proof}
Apply $\omega$: squares are nonnegative by positivity, constraints vanish
by admissibility, and $|\omega(R)| \le \|R\|$ by the C$^*$ norm bound.
\end{proof}

\begin{remark}[Relation to the moment--SOS and bootstrap literatures]
\label{rem:bootstrap}
The lower side of every emission is a noncommutative
moment--sum-of-squares bound, and the architecture (positivity of
moment matrices below, exhibited states above) is shared with active
programs: the variational reduced-density-matrix method
\cite{coleman,mazziotti}, the Navascu\'es--Pironio--Ac\'in hierarchy
\cite{npa}, and the Hamiltonian and lattice bootstrap
\cite{lin,hhk,andersonkruczenski,kazakovzheng}.  That line has advanced
well past its first papers: Euclidean lattice Yang--Mills has been
bootstrapped at finite $N$, including SU(2) in two, three, and four
dimensions with two-sided margins on the plaquette free energy at the
$0.1\%$ level against Monte Carlo \cite{kazakovzhengfinite}, with
SU(3) and abelian content following \cite{guosu3,lizhou}; matrix
quantum mechanics has reached high-precision ground-sector correlators
including D0-brane/BFSS content \cite{lindzero,linzhengmatrix}, whose
gravitational content is holographic (the bounded matrix theory is
dual to a black hole); gravitation enters the present paper in an
unrelated and more modest sense, as declared linear spin-2 content
with a stated exposure (\S\ref{sec:gravreads},
Remark~\ref{rem:deser}), and no holographic claim is made anywhere;
and
the quantum-information side has produced ground-energy lower bounds
and equilibrium algorithms with certified convergence guarantees
\cite{kull,fawzi}.  The present formulation is a member of this
family, and its worked instance (\S\ref{sec:instance}) is deliberately
the smallest nonsolvable case, not a competitive computation.  What it
adds is: the Hamiltonian (Kogut--Susskind) formulation posed directly
on the gauge-invariant coordinate ring, the region constraint being
the second fundamental theorem in person
(\S\ref{sec:correspondence}), where the cited lattice-gauge results
are Euclidean; the certification discipline, finite exact-rational
identities in the lineage of \cite{peyrlparrilo} with anytime validity
(Proposition~\ref{prop:weak}), where the cited programs report the
output of floating-point conic solvers; the two-sided rule that every
emission carry an exhibited admissible state on its upper side; the
relational emission doctrine; and the superstructure of
\S\S\ref{sec:canonical}--\ref{sec:continuum} (geometry, extensivity,
refinement) erected on the same cone.  The finite-$N$ Euclidean SU(2)
margins of \cite{kazakovzhengfinite} are the natural cross-community
contact point for the Hamiltonian brackets emitted here, once the two
formulations are related (anisotropic-limit matching).
\end{remark}

\begin{definition}[Boson word calculus]\label{def:bosoncalc}
When the content carries bosonic (Weyl/CCR) sectors the unbounded canonical
operators forbid a naive sum-of-squares reading, and certificate identities
are instead required to live in the $*$-algebra of boson-polynomial words
under the syntactic constraint that every term of positive boson degree
appear inside an exactly matched \emph{designated square} of boson degree at
most two per sector: a number form $b^\dagger b$ for any symplectic
(Bogoliubov) image $b = ua + va^\dagger$ of the budgeted family (the
original modes at $v=0$); a completed square
\[
-(Ba^\dagger + B^\dagger a) = \varepsilon\omega\,a^\dagger a
+ (\varepsilon\omega)^{-1}B^\dagger B
- (\sqrt{\varepsilon\omega}\,a + B/\sqrt{\varepsilon\omega})^\dagger
(\sqrt{\varepsilon\omega}\,a + B/\sqrt{\varepsilon\omega})
\]
with $B$ bounded; or a product of designated forms across commuting sectors
($[a,b]=0$ gives $(a^\dagger a)(b^\dagger b) = (ab)^\dagger(ab)$).  The
verifier checks the patterns exactly; residual words are boson-degree-zero,
and the anytime bound of Proposition~\ref{prop:weak} applies to them alone.
Admissibility carries the finite-moment proviso per sector, bounded through
the restricted covariance's spectrum.
\end{definition}

The physical state is selected linearly: $\w$ minimizes
$\omega \mapsto \omega(h)$ over admissible states.  The \emph{ground
sector} is the set of admissible states saturating the certified ground
energy; operationally it is the zero-temperature end of the completely
passive states \cite{pusz}.  For a nondegenerate gapped ground state the
sector is a point.

\begin{hypothesis}[Regularity, per content]\label{hyp:reg}
$h(\lambda)$ is essentially self-adjoint on the canonical core; the
ground state is unique on the systems considered; the quadratic form
$\E_\omega$ of \S\ref{sec:dirichlet} is closable.  These are standard and
are verified per instance; all results below assume them where used.
\end{hypothesis}

\subsection{The upper bracket: the trial ladder}\label{sec:ladder}
Weak duality supplies the lower side of every emission; the upper side is
the exact evaluation of an exhibited admissible state, and it is as
constructive as the lower.  An exhibited state meets four requirements, the
trace of the tower's own first level: (T1) it is admissible and saturates
the ground energy $e_0$, or averages to such a state at unchanged value
under the content's compact symmetry; (T2) its per-covolume moments of
local words are finite exact contractions with certified rounding
(uncertified tails are forbidden, tails as such are not); (T3) one
construction serves every content; (T4) no object outside the theory
enters, the class being generated from a reference determined by the
content by operations internal to $\A$.  The reference is one object, read
at the content's group.

\begin{proposition}[The reference rung is the kinetic vacuum]
\label{prop:reference}
The reference state is the ground state of the kinetic Hamiltonian
$h(0) = \mathrm{Cas}$: the invariant vacuum of the content's symmetry,
determined by the content alone with no choice, whose moment functional is
that symmetry's singlet-channel fusion: the value of
$\omega_{\mathrm{ref}}$ on a product of words is the trivial-representation
content of the corresponding tensor product.  For compact gauge content
this is the Haar state, its moments computed by character fusion
$\int_G \chi_{j_1}\cdots\chi_{j_n}\,dU = \dim\mathrm{Inv}(V_{j_1}\otimes
\cdots\otimes V_{j_n})$ (\S\ref{sec:instance}); for CAR$\otimes$CCR field
content it is the quasi-free$\otimes$Gaussian state, its moments computed by
Wick contraction (\S\ref{sec:subsumption}).  These are one construction at
two groups: the pair contraction is the Heisenberg (or Clifford) group's
singlet projector on $V\otimes V$ exactly as character orthogonality is the
compact group's, so Wick's theorem and character fusion are the same
operation: summing the singlet channels of a tensor product.
\end{proposition}

\begin{remark}[The compact seam]\label{rem:refseam}
For a compact group the invariant vacuum is the tracial Haar state; for a
noncompact field sector it is the Fock vacuum of the distinguished kinetic
representation.  The difference is in what the vacuum \emph{is} (a trace
versus a lowest vector), not in the construction: both are the $h(0)$
ground state and both carry the singlet-fusion moment theorem.  This is the
sole place the one construction touches compactness; it is not a second
construction.  \S\ref{sec:correspondence} gives the seam structural
form: compact and field content are the two members of a reductive dual
pair on one carrier, and the compression between them is the Gauss law.
\end{remark}

\begin{definition}[Enrichment class and trial ladder]\label{def:ladder}
The enrichment class $\mathfrak E$ is the automorphisms of $\A$ commuting
with the content's symmetry and \emph{strictly local} (each local element
maps to a local element within a finite radius), graded by radius;
concretely they are generated by the strictly-local elements from which the
magnetic element $M$ is built: plaquette (Wilson-loop) operators for
gauge content, finite-radius charge-conserving circuits for field content,
one operation with $M$ the only input.  Rung $r$ of the trial ladder is the
orbit of the reference rung under radius-$\le r$ elements of $\mathfrak E$,
conditioned by the character integrals of any sharp sectors; $U^*(r)$ is
the ground-sector infimum of the read over rung $r$, non-invariant members
averaged.  Moments are exactly evaluable at every rung: conjugation of a
local word through a radius-$r$ element is exactly finite (a Pfaffian, a
character fusion, or a Gaussian characteristic function, closed for
symplectic and mixed-even gates, norm-convergent with absorbed tails
beyond) at cost $e^{O(r)}\,\mathrm{poly}$ in the lightcone volume.
\end{definition}

\begin{proposition}[Determination]\label{prop:ladder}
The trial ladder is the unique class satisfying (T1)--(T4) with intrinsic
richness grading: the reference is forced by
Proposition~\ref{prop:reference}, the operations by the trace of
(T1)--(T4), the grading is the enrichment radius.  Each familiar
alternative fails a named requirement.  A similarity-transformed
expectation that is not a state evaluation emits no bound (T1); a
non-truncating unitary transform is not exactly evaluable (T2); a
stochastic estimator carries error bars, not certificates (T2); a tensor
network exactly contractible only in one dimension fails (T3), and
elsewhere its approximate contraction bounds nothing (T2); a
density-functional model is a fitted surrogate emitting no bound (T1, T2).
The licensing runs forward: any exhibited admissible ground-sector state
with exactly evaluable moments emits a valid bound, on or off the ladder.
\end{proposition}

\subsection{Residual nonlinearity, classified}\label{sec:nonlinearity}
Exactly two kinds of nonlinearity in the state survive audit, and both
are linear structure in disguise: (i) \emph{envelopes}: $E_1(\lambda) =
\inf_\omega[\omega(\mathrm{Cas}) - \lambda\,\omega(M)]$ is a lower
envelope of affine functions of $\omega$; its concavity in $\lambda$ and
the Hellmann--Feynman relation $dE_1/d\lambda = -\langle M\rangle$ are
the supporting-hyperplane structure of the moment cone; (ii)
\emph{content algebra}: polynomial relations among canonical generators
(e.g.\ the cubic relation of \S\ref{sec:instance}) are algebraic geometry
of the content, not of the state.  Anything else nonlinear (wavefunction
representatives, field thresholds, rescaled block variables) is a chart
and is excluded from definitions and emissions.

\subsection{The uncertainty hierarchy}\label{sec:uncertainty}
Membership in the moment cone is a hierarchy of uncertainty relations, and
the width of every emission is a read of the state's distance from its
boundary.  Fix the degree-$\le n$ word space $\W_n =
\mathrm{span}\{w_\alpha\}$ ($\dagger$-closed, $w_0 = \mathbf 1$) and the
level-$n$ moment matrix $M_n(\omega)_{\alpha\beta} = \omega(w_\alpha^\dagger
w_\beta)$.

\begin{theorem}[Generalized uncertainty at level $n$]\label{thm:uncertainty}
Write the centered Gram $G^c_{\alpha\beta} = \omega(\delta w_\alpha^\dagger
\,\delta w_\beta)$, $\delta w_\alpha = w_\alpha - \omega(w_\alpha)$, and
split it into Hermitian symmetric and commutator parts $G^c = \Sigma +
\tfrac12 K$ with $\Sigma_{\alpha\beta} = \tfrac12
\omega(\{\delta w_\alpha^\dagger, \delta w_\beta\})$, $K_{\alpha\beta} =
\omega([\delta w_\alpha^\dagger, \delta w_\beta])$.  Then level-$n$
admissibility $M_n \succeq 0$ is exactly $\Sigma \succeq -\tfrac12 K$, and
straightening the canonical relations evaluates $K$ as an affine function
of moments of degree at most $2n-2$.  At $n=1$ on a bosonic sector $K$ is
the central symplectic form and this is the Robertson--Schr\"odinger
relation; at $n=1$ on a fermionic sector the anticommutator matrix is
central and $M_1 \succeq 0$ is the Coleman constraint $0 \preceq \gamma
\preceq \mathbf 1$.  Heisenberg uncertainty and the exclusion principle are
the two graded instances of one statement, and level $n$ is its degree-$n$
extension.
\end{theorem}
\begin{proof}
$w_\alpha^\dagger w_\beta = \tfrac12\{w_\alpha^\dagger, w_\beta\} + \tfrac12
[w_\alpha^\dagger, w_\beta]$ gives $G^c = \Sigma + \tfrac12 K$ after
centering (a Schur complement in the unit entry); $\Sigma$ and $K$ are
Hermitian, so $G^c \succeq 0 \iff \Sigma \succeq -\tfrac12 K$, identified
with $M_n \succeq 0$ by centering.  The commutator of degree-$d_\alpha,
d_\beta$ words straightens to degree $\le d_\alpha + d_\beta - 2 \le 2n-2$.
\end{proof}

\begin{theorem}[Gap identity]\label{thm:gapid}
For a level-$n$ certificate with Gram $\Lambda = \sum_i c_i f_i f_i^\dagger
\succeq 0$ (factors $F_i = \sum_\alpha (f_i)_\alpha w_\alpha \in \W_n$) and
emitted value $\lambda$, and any admissible ground-sector state $\omega_0$
with carrier vector $\Omega_0$,
\[
\omega_0(X) - \lambda = \sum_i c_i \|\pi(F_i)\Omega_0\|^2
= \langle M_n(\omega_0),\, \Lambda\rangle_F .
\]
The shortfall of any certificate on the physical state is the
$\Lambda$-weighted squared residual of the state under the certificate
factors: the squared distance, \emph{in the certificate's own seminorm}, of
the state from the kernel variety $\{F_i\psi = 0\}$.  The emitted width is
this distance, not a generic numerical residue.
\end{theorem}
\begin{proof}
Evaluate the certificate identity $X - \lambda\mathbf 1 - (\text{data
multipliers}) = \sum_i c_i F_i^\dagger F_i$ in $\omega_0$: the energy and
data multipliers vanish on the admissible ground sector, and
$\omega_0(F_i^\dagger F_i) = \|\pi(F_i)\Omega_0\|^2 = f_i^\dagger
M_n(\omega_0) f_i$ sums to the Frobenius pairing.
\end{proof}

\begin{theorem}[Closure is boundary contact]\label{thm:closure}
The bracket closes at level $n$, $L_n^* = G^*$, iff the physical moment
matrix sits on the face of the level-$n$ cone exposed by the optimal
certificate: for an attained optimum $(\lambda^*, \Lambda^*)$, the vanishing
pairing $\langle M_n(\omega_0^*), \Lambda^*\rangle_F = 0$ with both factors
positive semidefinite forces $M_n(\omega_0^*)\Lambda^* = 0$, so every
certificate factor annihilates the state.  This generalizes minimum
uncertainty (degree-one annihilation) to a degree-$n$ hierarchy: the
degree-one kernel states are the squeezed and coherent vacua; the
degree-two boundary already contains the Fock states ($(a^\dagger a -
m)|m\rangle = 0$), all squeezed states, and the even/odd cat states ($a^2 -
\alpha^2$ annihilates the even cat), on which level-two brackets close
exactly.  A ground state annihilated by no degree-$\le n$ system compatible
with the multipliers has positive width at every finite certificate, and
the tower's convergence is the statement that these residuals fall to zero
as $n$ grows.
\end{theorem}
\begin{proof}
$\Tr(AB) = 0$ for $A, B \succeq 0$ gives $AB = 0$; applied to
Theorem~\ref{thm:gapid} at the extremal state, $f_i^\dagger M_n f_i = 0$ for
each $c_i > 0$, i.e.\ $\|\pi(F_i)\Omega_0^*\| = 0$.  The converse is the gap
identity read on an annihilating certificate; without attainment the
statement holds in the supremum, and every finite certificate still emits a
valid bound.
\end{proof}

\begin{remark}[The width is not a local defect]\label{rem:widthdefect}
A local uncertainty defect is not the emitted width.  A two-mode Gaussian
ground state whose global normal-mode annihilators are the certificate
factors has certified width $4.9\times10^{-6}$ at its stated budget (the
bracket closes as the budget grows), while the
Robertson--Schr\"odinger defect of its reduced single mode is $\nu^2 -
\tfrac14 = 0.1406$: Theorem~\ref{thm:gapid} governs, and its distance is
measured in the certificate seminorm, whose active directions the
multipliers select, never in a fixed local metric.
\end{remark}

The annihilating factors of Theorem~\ref{thm:closure} are exactly the
degree-$\le n$ elements of the Gelfand ideal $N_\omega$ of
\S\ref{sec:canonical}: exactness is ideal membership.  In the
Hamiltonian-bootstrap literature the same phenomenon appears as the
null-state condition \cite{guoli}.  The quantitative
converse (a lower bound on the width by a fixed-metric distance to the
exposed faces, with constants uniform over reads and budgets) is an
error-bound question for the noncommutative hierarchy, and is not claimed.

\subsection{Resolving emissions}\label{sec:resolving}
A bracket that contains a target value confirms nothing by itself.  An
emission \emph{resolves} a comparison only when its certified width is
below a threshold declared in advance of the computation; thresholds may
be tightened before a computation and never loosened after one.  All
comparisons in this paper are subject to this rule.

\subsection{Execution and cost}\label{sec:execution}
The instrument is a constructor, not an existence argument: every emission
is produced, and its execution has exactly four kinds of step, none of them
a self-consistency loop.  (1) One direct-optimization stream per
certificate pullback, a search on a merit function whose convergence is not
load-bearing (Proposition~\ref{prop:weak}).  (2) Machine-precision
primitives that terminate: Gauss--Hermite integrals, characteristic
functions and Pfaffians, character fusions, linear solves.  (3) A finite
verification pass in exact or interval arithmetic, matched patterns
(Definition~\ref{def:bosoncalc}) checked syntactically.  (4) Monotone budget
sweeps with warm starts.  The certificate multiplier is carried as a Gram
$\Lambda = FF^\dagger$ with the factors transient; the scalar multipliers
restore by closed-form least squares given $F$; the climb maximizes the
verified emitted value, unconstrained over $F$, with rank a monotone budget
axis.  That no step is a fixed-point iteration is what makes soundness
independent of convergence: a valid two-sided bracket is in hand at every
step of every stream, and the streams only narrow it.

\begin{theorem}[Per-sweep optimality]\label{thm:persweep}
With $M$ the generator-family size per covolume, the level-$n$ certificate
identity for a fixed matrix element has $\Theta(M^{2n}/(2n)!)$
independent coefficients up to the content's word-symmetry factor (a
constant of the content), each verified once, and the climb's $O(M^{2n})$
per sweep is input-linear in that count, optimal up to rank and constants.
A conditional projection onto a commuting subalgebra acts within the fixed
word family and does not change the count; the energy-saturation multiplier
adds one scalar.  The counting is content-blind: each specialization of a
content preserves the scaling class of the one it refines, paying an
additive family count.  Total time is $I(\epsilon)\cdot O(M^{2n})$, with
$I(\epsilon)$ the number of sweeps spent to reach emitted width
$\epsilon$, a measured quantity carrying no a priori bound, so no
unconditional total-time claim is made; natively sum-of-squares pieces
(reciprocal kernels, matched boson squares) verify below the generic
floor.
\end{theorem}

\section{Canonical structures of the state}\label{sec:canonical}

Throughout, $\w$ is the ground state of declared content under
Hypothesis~\ref{hyp:reg}.

\begin{remark}[Convention on canonicity]\label{rem:convention}
``Canonical'' below always means: determined by the content and the
state up to an equivalence that is named at the construction and that
acts trivially on every emission.  Emissions are values of the state on
algebra elements; unitary equivalence of the carrier and the inner
ambiguity of modular data leave every such value fixed, so the residual
equivalence classes are gauge that never reaches an emission.  Wherever
a construction below depends on more than $(\text{content}, \w)$ modulo
such an equivalence, that dependence is stated, not absorbed into the
word.
\end{remark}

\subsection{Carrier and closure}
The GNS construction gives the carrier $(\pi, \Hw, \Omega)$, unique up to
unitary equivalence.  The \emph{Gelfand ideal} $N_\omega = \{a \in \A :
\omega(a^\dagger a) = 0\}$ is determined by the state alone.  By
\S\ref{sec:uncertainty}, $N_\omega$ is the closure structure of the
instrument: brackets close exactly on certificates whose factors lie in
it.

\begin{hypothesis}[Faithfulness]\label{hyp:faithful}
$\w$ is faithful on the abelian (multiplication) subalgebra.  This is the
chart-free residue of the classical statement ``the ground wavefunction
is strictly positive''; it is verified on the instances considered.
\end{hypothesis}

\subsection{Modular structure}\label{sec:modular}
For a subalgebra $\mathcal N \subseteq \pi(\A)''$ on which $\Omega$ is
cyclic and separating, Tomita--Takesaki theory yields the modular data
$(\Delta_{\mathcal N}, J_{\mathcal N})$ and flow, canonical up to inner
automorphisms \cite{takesaki}.  A pure ground state is separating for no
global type-I algebra; modular structure therefore lives on
\emph{subalgebras}, and the admissibility criterion is internal:
$\mathcal N$ is admissible precisely when $\mathcal N \cap N_\omega =
\{0\}$.  Time, in this formulation, is the modular flow of admissible
subalgebra--state pairs; it is an output.

\subsection{Dirichlet form, Dirac operator, factorization}
\label{sec:dirichlet}
The canonical derivation is $\partial f = (X f)\,\Omega \in \Hw \otimes
\g$, valued in the Clifford module of $(\g, \mathrm{Killing})$; the
\emph{Dirichlet form of the state} is
\[
\E_\omega(f,g) \;=\; \sum_a \omega\big((X_a f)^{*}(X_a g)\big),
\]
manifestly linear in $\omega$.  Its generator structure is the
Cipriani--Sauvageot theory \cite{cipriani}; the choice of first-order
differential structure on which that theory depends is not made
downstream: the derivations are declared content (\S\ref{sec:content}),
so the construction is functorial in content, and its content dependence
is the dependence every object of the theory has.  The associated Dirac
operator $D_\omega$ satisfies $D_\omega \Omega = 0$ by construction.

\begin{theorem}[Harmonicity]\label{thm:harmonicity}
Let the content be of Dirichlet type (the class of
Theorem~\ref{thm:factorization} below) and assume
Hypotheses~\ref{hyp:reg}--\ref{hyp:faithful}.  Let $\omega$ be an
admissible state whose dynamics is implemented on its carrier: a
self-adjoint $H_\omega$ with $\Omega$ in its form domain and
$e^{itH_\omega}\pi(a)e^{-itH_\omega} = \pi(\alpha_t(a))$ on the
canonical core.  Then the identity
$\E_\omega(f,g) = \langle f\Omega, (H_\omega - E_0)\, g\Omega\rangle$
holds on the core if and only if $\omega$ is the ground state.  The
diagonal defect at $f = g = \mathbf 1$ is $\omega(h) - E_0$: a
nonnegative affine read of the state (the \emph{geometricity defect}
at unit), vanishing exactly on the ground sector.
\end{theorem}
\begin{proof}
If $\omega$ is the ground state, the identity is
Theorem~\ref{thm:factorization}.  Conversely, evaluate the claimed
identity at $f = g = \mathbf 1$: the left side vanishes because
$X_a\mathbf 1 = 0$, and the right side is
$\langle\Omega, (H_\omega - E_0)\Omega\rangle = \omega(h) - E_0$, so
$\omega(h) = E_0$ and $\omega$ lies in the ground sector; uniqueness
(Hypothesis~\ref{hyp:reg}) identifies it as the ground state.
\end{proof}

\begin{theorem}[Automatic local factorization, stoquastic class]
\label{thm:factorization}
Let $h = \mathrm{Cas} - \lambda M$ be of Dirichlet type with respect to
the declared derivations: in the content's function realization, $h$ is
an elliptic Schr\"odinger operator (Casimir Laplacians per link plus a
real multiplication potential), so that the ground representative
$\psi_0$ is strictly positive (Perron--Frobenius; this is
Hypothesis~\ref{hyp:faithful} holding by theorem rather than by
assumption, and pure gauge content is in the class).  Then, in the
ground carrier,
\[
\langle f\Omega,\,(h - E_0)\,g\Omega\rangle
= \sum_a \omega\big((X_a f)^*(X_a g)\big),
\qquad\text{i.e.}\qquad
h - E_0 = \sum_a L_a^{\dagger} L_a,
\]
with $L_a(f\Omega) = (X_a f)\Omega$ local to link $a$ and $L_a \Omega =
0$ for every $a$.
\end{theorem}
\begin{proof}
On the class stated, $e^{-th}$ is positivity improving and $\psi_0 > 0$
(Perron--Frobenius for Schr\"odinger operators on products of compact
group manifolds).  The ground-state (Doob) transformation, standard
in Dirichlet-form theory \cite{fot}, then gives
$(J_a - i(X_a\phi))(g\psi_0) = -i(X_a g)\psi_0$ for $\phi = -\log
\psi_0$, by the Leibniz rule; summing adjoint squares yields the
displayed identity, whose two sides are chart-free and linear in
$\omega$, hence hold in the carrier independently of the realization.
Locality holds because $X_a$ acts on the tensor factor of link $a$;
$L_a\Omega = 0$ because $X_a 1 = 0$.  Domains: the identity holds on the
canonical core, where both sides are densely defined closable quadratic
forms; Hypothesis~\ref{hyp:reg} (essential self-adjointness and
closability) identifies their closures, extending the identity to the
form domains.
\end{proof}

\begin{remark}[Scope of the automatic precondition]\label{rem:scope}
On the stoquastic class of Theorem~\ref{thm:factorization}, the ground
carrier automatically presents the Hamiltonian as a sum of local
positive operators each annihilating the ground vector: the structural
precondition of the finite-size gap criteria of
\S\ref{sec:extensivity}, which generic Hamiltonians fail.  The scope
boundary is exactly the sign problem.  Fermionic hopping content is not
of Dirichlet type with respect to the declared derivations; the ground
representative of such content is not positive in any local realization
in general, and the obstruction cannot be removed by local basis
changes; curing non-stoquasticity is computationally hard even when
possible \cite{bdot,mlh}.  For that content the transformation above is
unavailable, and the identity between the ground form and the content
Dirichlet form is not expected to hold.
\end{remark}

\begin{openproblem}[Local factorization for fermionic content]
\label{op:sign}
Determine for which fermionic contents the ground form admits any local
positive decomposition $h - E_0 = \sum_x Q_x$ with $\Omega \in \ker Q_x$
for every $x$.  This is the sign problem in carrier form.  It is not an
item completed by additional work of a known kind: a structural
obstruction is possible, and either an existence argument or an
obstruction theorem would be a result.  Absent a decomposition, the
certified finite-size gap criterion of \S\ref{sec:extensivity} does not
apply to fermionic content at moderate coupling, and the gap program for
such content routes through the strong-coupling regime alone
(\S\ref{sec:gaps}).  A candidate structural criterion is
Conjecture~\ref{conj:signfree} (\S\ref{sec:correspondence}).
\end{openproblem}

\subsection{Metric, dimension, volume, sectors}\label{sec:metricdim}
The carr\'e du champ $\Gamma(f,g) = \tfrac12[f\,\mathfrak L g + g\,
\mathfrak L f - \mathfrak L(fg)]$, with $\mathfrak L$ the Casimir
Laplacian, is content-canonical.  Points are pure states of the abelian
sector; the distance
\[
d(p,q) = \sup\{\,p(f) - q(f) : \Gamma(f,f) \preceq 1\,\}
\]
is a conic linear program in the instrument's own certificate cone
\cite{connes}.  The counting function $N(E)$ of $D_\omega^2 = h - E_0$
is a linear read of spectral projections; spectral dimension is a
certified window slope of $N$; spectral volume is the Weyl coefficient,
admissible as an emission only where Dixmier measurability (existence of
the limit) is itself certified.  Quantum-number sectors are the gradings
under the automorphism group of the content; all emissions are invariant
under them.  On instances with unique ground state, $\dim\ker D_\omega =
1$.

\section{The geometric identification}\label{sec:geometry}

The identification made by this formulation is: \emph{the certificate
structure of the state, at closure, is a Dirac geometry}.  The data
$(\A, \Hw, D_\omega)$ constitute a spectral triple only if the axioms
hold: compact resolvent of $D_\omega$, bounded commutators $[D_\omega,
f]$ on a dense subalgebra, regularity, and finite dimension spectrum
\cite{connes}; a manifold requires in addition the reconstruction axioms
\cite{connesrecon}.  None of these is assumed here.  On finite content
they are finite-dimensional trivialities (bounded commutators are the
finiteness of $\Gamma(f,f)$ already used by the metric), and the
substantive questions are their uniformity along extensivity
(\S\ref{sec:extensivity}) and refinement (\S\ref{sec:continuum}); each
axiom is a named verification obligation per instance, on the same
footing as Dixmier measurability in \S\ref{sec:metricdim}.  The
identification is falsifiable at its first quantitative joint:

\begin{definition}[Weyl consistency requirement]\label{def:weyl}
For every instance, the certified spectral volume (the Weyl coefficient
of $N(E)$, where measurable) must equal the Hausdorff volume of the
spectrum under the metric $d$.  Failure of this equality at certified
resolution refutes the identification for the instance.
\end{definition}

For the worked instance of \S\ref{sec:instance} both sides are in
closed form.  The metric side (all values derived in
Appendix~\ref{app:geometry}): the required volume is
$(32 - 16\sqrt 2)\pi^2$, the diameter
realization $d(\mathbf 1, (-\mathbf 1,-\mathbf 1)) = 2\sqrt2\,\pi$, and
the one-loop distance $2\,|\arccos(x_1/2) - \arccos(x_2/2)|$.  The
spectral side is also closed form for this instance: every eigenvalue
and multiplicity of $D_\omega^2$ on the carrier is exact
(Appendix~\ref{app:weyl}), so the eigenvalue count is exact arithmetic,
the existence of the Weyl limit is proved rather than assumed (Dixmier
measurability holds), and the comparison closes at every threshold of
the resolution rule of \S\ref{sec:resolving}.  The comparison has been
run, and it passes exactly:

\begin{theorem}[Weyl consistency holds for the worked instance]
\label{thm:weylinstance}
For the content of \S\ref{sec:instance} the counting function of
$D_\omega^2 = h - E_0$ satisfies
\[
N(E) \;=\; \frac{16 - 8\sqrt2}{3}\,E^{3/2} \;+\; O(E),
\]
and the Weyl coefficient equals the Hausdorff volume of
Appendix~\ref{app:volume} as an algebraic identity:
\[
6\pi^2 \cdot \frac{16 - 8\sqrt2}{3} \;=\; (32 - 16\sqrt2)\,\pi^2 .
\]
The spectral dimension read is $3$, the dimension of the metric side.
The leading coefficient is coupling-independent: $\|\lambda M\| \le
6\lambda$ sandwiches $N_{h-E_0}$ between shifted copies of the kinetic
counting function, so the identity holds at every declared coupling.
Proof in Appendix~\ref{app:weyl}.
\end{theorem}

\noindent The two closed forms share no machinery: the volume is a
triple integral of the $\Gamma$-metric's volume element, the count is
Clebsch--Gordan combinatorics; their agreement is exactly the equality
Definition~\ref{def:weyl} was built to test, and it is an instance of
the equivariant Weyl law \cite{donnelly,bruningheintze}, verified
directly.  On a commutative carrier
the identification \emph{must} reduce to the classical Weyl law, so
the content of the test is not the law but the normalizations: every
constant entering either side was fixed elsewhere (the metric
normalization by the verified carr\'e du champ of
\S\ref{sec:instance}, the operator by the declared content, the
quotient measure by the submersion), with no adjustable element
remaining at the comparison.  A mismatch by any factor would have
refuted the identification; there was none.  No result elsewhere in
the paper depends on this outcome; the geometric identification has
now passed its first decisive test, and the continuum program of
\S\ref{sec:continuum} stands or falls independently.

\section{Worked instance: SU(2) on two plaquettes}\label{sec:instance}

\subsection{Content and algebra}
Gauge group SU(2) on two faces sharing a boundary; the invariant algebra
is generated by the loop characters $x = \chi(U_1)$, $y = \chi(U_2)$,
$z = \chi(U_1 U_2)$, freely up to the classical relation defining the
character variety
\[
\kappa(x,y,z) \;=\; x^2 + y^2 + z^2 - xyz - 4 \;\le\; 0,
\qquad |x|,|y| \le 2,
\]
and the Hamiltonian is $h(\lambda) = C_2^{(1)} + C_2^{(2)} -
\lambda(x+y+z)$.  All moments of the Haar state are exact rationals
computed by character fusion; the carr\'e du champ closes on polynomials
in $(x,y,z)$ with rational structure constants (e.g.\ $C_2\,x =
\tfrac34 x$, $\Gamma(x,x) = (4 - x^2)/4$).

\subsection{Certified spectra}
Two-sided brackets on $E_1(\lambda)$, produced by sum-of-squares
certificates with multipliers on the region above (lower side) and exact
evaluation of exhibited states (upper side), verified in exact rational
arithmetic:

\begin{center}
\begin{tabular}{lll}
\toprule
$\lambda$ & certified bracket on $E_1$ & width \\
\midrule
$1/16$, $1/12$ & two-sided, width $\le 10^{-15}$ & closed \\
$1/2$ & $[-0.9548,\ -0.8836]$ & $0.071$ \\
$2$ & $[-7.0500,\ -6.9595]$ & $0.090$ \\
\bottomrule
\end{tabular}
\end{center}

\noindent The plaquette expectation obeys the certified bounds
$4.0031 \le \langle x{+}y{+}z\rangle_{\lambda=2} \le 6$, the lower side
from the concavity (supporting-hyperplane) structure of
\S\ref{sec:nonlinearity}, the upper from the exact norm.

\subsection{Coupling towers}
\emph{Weak coupling} (derived; Appendix~\ref{app:fourth}):
\[
E_1(\lambda) = -\tfrac{10}{3}\lambda^2
- \tfrac{32}{9}\lambda^3 + \tfrac{551}{297}\lambda^4
+ \tfrac{172064}{9801}\lambda^5 + O(\lambda^6),
\]
with all coefficients exact and consistent with the certified
brackets.  \emph{Strong coupling}
(derived at leading orders; Appendix~\ref{app:strong}): harmonic
expansion about the identity class in the $\Gamma$-metric gives
normal-mode stiffnesses $3\lambda/4$ and $\lambda/4$ (three components
each), hence
\[
E_1(\lambda) = -6\lambda + \tfrac32(1{+}\sqrt3)\sqrt\lambda + c_0 +
O(\lambda^{-1/2}),
\qquad
E_2 - E_1 = 2\sqrt\lambda + O(1),
\]
gap-ratio asymptote $(E_3 - E_1)/(E_2 - E_1) \to (1+\sqrt3)/2$, plaquette
law $\langle x{+}y{+}z\rangle = 6 - \tfrac34(1{+}\sqrt3)/\sqrt\lambda +
O(1/\lambda)$ (consistent with Hellmann--Feynman by construction and with
the certified bounds above); the invariant-sector selection rules place
the swap-even excitation lowest at both ends of the coupling axis.  The
constant term is derived in Appendix~\ref{app:strong}:
\[
c_0 \;=\; -\frac{126 + 13\sqrt3}{192} \;\approx\; -0.77352,
\]
inside the certified window $-0.79 \pm 0.05$ at $\lambda = 2$
outright; fixing it determines subleading coefficients
$(c_{-1}, c_{-2}) \approx (+0.07, -0.08)$ from converged variational
values at $\lambda \in \{2, 8\}$, which then land the certified window
$-0.83 \pm 0.04$ at $\lambda = 1/2$---one constant and two small
corrections carrying four couplings.

\section{The field--gauge correspondence: the group as output}
\label{sec:correspondence}

Proposition~\ref{prop:reference} exhibited one construction at two
groups.  This section states the structure behind the seam of
Remark~\ref{rem:refseam} and draws its consequence: compact-gauge
content and CAR$\otimes$CCR field content are two generating views of a
single carrier, related by a reductive dual pair; the Gauss law is the
singlet compression between the views; and, on this correspondence, the
symmetry itself moves from the input side to the output side.  The
group's rank is recovered from the closure ideal of the reference
state, the same object that closes brackets
(Theorem~\ref{thm:closure}), and the existence of a symmetry becomes a
positivity question of the type the instrument already runs.

\subsection{One carrier, two generating subalgebras}
\begin{definition}[Dual-pair carrier]\label{def:dualpaircarrier}
Let $W = W_1 \oplus W_0$ be a $\mathbb Z_2$-graded one-particle space
and $\mathcal F$ the Fock space carrying CAR over $W_1$ and CCR over
$W_0$.  The quadratic words close on the orthosymplectic algebra
$\mathfrak{osp}(W)$, the kinetic symmetry of field content (the
Bogoliubov transformations).  A \emph{reductive dual pair} $(\g, \g')$
is a pair of subalgebras of $\mathfrak{osp}(W)$, each the full
commutant of the other \cite{howe}.  The distinguished vector is the
lowest-weight vector $\Omega$ of the oscillator representation, fixed
by the compact stabilizer that the kinetic Hamiltonian's spectral split
determines; as in \S\ref{sec:sea-instance}, the polarization is an
output of the content, not a choice, and no fully
$\mathfrak{osp}$-invariant state exists to choose instead (for CCR none
exists; for CAR it is the trace, not the Fock vacuum).
\end{definition}

The \emph{gauge view} generates from $\Omega$ by the $\g$-invariant
words (loop words and characters); the \emph{field view} generates by
the linear field words.  The kinetic term of the gauge view is the
electric Casimir of $\g$; that of the field view is the quadratic
current form $\sum_a J_a J_a$ of the dual partner; the two Casimirs of
a dual pair are related through the number operator by Capelli-type
identities.  Both views read one state.

\begin{proposition}[Wick is fusion; the Gauss law is Haar integration]
\label{prop:wickfusion}
The moments of $\Omega$ on $\g$-invariant words equal the
singlet-channel fusion of the corresponding tensor products: both sides
evaluate the projector onto the trivial isotypic component.  The first
fundamental theorem of invariant theory supplies the dictionary
(pairings span the invariants of tensor powers: symmetric for
orthogonal, antisymmetric for symplectic, mixed $V$--$\bar V$ for
unitary content) \cite{weylbook}, and the Weingarten calculus is its
exact finite-$N$ form: Haar integration of a word equals a signed sum
of Wick pairings with explicit rational coefficients
\cite{collinssniady}.  Consequently the compression
$P = \int_G \pi(U)\, dU$ onto the singlet sector is Haar integration,
is linear in the state, and forms no section: the gauge state is the
field state compressed to invariant words, in the discipline of
Remark~\ref{rem:sections}.  Provenance: classical invariant theory
transported whole, numbered for the dependency graph
(\S\ref{sec:conclusion}).
\end{proposition}

\subsection{Exactness and budget}
\begin{proposition}[Prepotential exactness for SU(2)]\label{prop:prepot}
Peter--Weyl gives $L^2(\mathrm{SU}(2)) \cong \bigoplus_j V_j \otimes
V_j$, the number-matched singlet sector of the Fock space of two
oscillator doublets per link, one per link end \cite{schwinger,mathur}.
Under the identification the electric Casimir is polynomial in the
number operator, $C_2 = (\hat N/2)(\hat N/2 + 1)$, and the loop
characters are polynomial words in oscillator bilinears.  The unitary
link operator itself carries $(\hat N + 1)^{-1/2}$ normalizations and
is rational, not polynomial; the instrument's generating set is the
polynomial intertwiner words, and the classical relation of
\S\ref{sec:instance} is carried by the closure ideal.  The
identification is exact: the worked instance of \S\ref{sec:instance}
is a singlet-compressed bosonic content with no truncation.
\end{proposition}

For general compact $G$ the known exact presentations are per
representation: rishon and quantum-link content realizes the link
algebra on a finite-dimensional fermionic carrier exactly at fixed
representation \cite{quantumlink}, so the representation cutoff is a
\emph{declared budget}, and full $L^2(G)$ is its refinement limit along
the machinery of \S\ref{sec:continuum}, a new budget axis with the same
exact projection structure (Proposition~\ref{prop:projection}).

\begin{openproblem}[The general dual-pair presentation]
\label{op:dualpair}
Construct, for each compact $G$ of the declared contents, the dual-pair
carrier and singlet compression presenting $L^2(G)$-content exactly
(SU(2) is exact above; SU(3) prepotentials with additional constraints
are known \cite{anishetty}), and carry the canonical machinery of
\S\ref{sec:canonical} through the compression.  This construction
routes Open Problem~\ref{op:car} and the gauge instance through one
object: the Dirichlet form, derivations, and $\Gamma$-calculus of the
gauge sector become Wick-computable, and the four tests of
Hypothesis~\ref{hyp:gma} and the gauge machinery stop being two
machineries.
\end{openproblem}

\subsection{The group is an output of the state}
\begin{proposition}[Group from closure]\label{prop:groupclosure}
For the classical families ($\mathrm U(N)$, $\mathrm O(N)$,
$\mathrm{Sp}(N)$) the moment matrix of the reference functional on
invariant words of degree $\le b$ is the Gram matrix of the pairing
(Brauer or partition-algebra) basis, with the rank $N$ entering only as
a parameter; the first fundamental theorem says these words span
\cite{weylbook}.  The family is fixed by the declared pairing symmetry
of the words (orthogonal, symplectic, or unitary), and the rank is
recovered from the Gelfand ideal of the functional: at integer $N$ the
ideal is generated by the relations of the second fundamental theorem,
the determinantal and Pfaffian identities, entering at a degree that
reads off $N$.  Content may therefore declare the fusion functional,
and the symmetry is read from the state by the instrument's own closure
test (\S\ref{sec:uncertainty}) rather than declared.  For the worked
instance the recovered relation is the Fricke cubic of
\S\ref{sec:instance}: the region constraint the instrument has been
carrying is the second fundamental theorem in person, its inequality
the real (positivity) locus.
\end{proposition}

\begin{remark}[What the ideal cannot see]\label{rem:complexpoint}
The output of Proposition~\ref{prop:groupclosure} is the
\emph{complexification} and nothing finer.  Invariant theory sees the
rank and the symmetry type of the invariant form, never its signature:
the Brauer data at $\delta = 4$ with symmetric pairing carry the same
fundamental theorems, hence the same closure ideal, for $O(4)$,
$O(3,1)$, and $O(2,2)$: the real forms of the one complex group
$O(4,\mathbb C)$ that preserve a real symmetric pairing.  (The
remaining real form, the quaternionic $O^*(4)$, preserves no real
symmetric form and corresponds to no signature; it is addressed with
the signature question in \S\ref{sec:realform}.)  This is a scope statement, not a defect: real forms
are classified by conjugacy classes of antilinear involutions, which no
complex-linear moment datum can encode, and the theory's canonical
antilinear objects are the modular conjugations of \S\ref{sec:modular}.
The resolution of the real form is therefore a modular question and is
taken up in \S\ref{sec:realform}, in that section's conditional
register.  Two finer points share this boundary: the ideal fixes $O$
rather than $SO$, the orientation word (the $\varepsilon$-tensor, a
Pfaffian invariant) not being in the pairing image, so orientation is a
separate declared-or-read item; and the integer output is the parameter
$\delta$, identified with a dimension only through the read of
\S\ref{sec:metricdim}.
\end{remark}

\begin{remark}[Quantization of symmetry as positivity]\label{rem:jones}
Away from the integer points the fusion functional on the pairing
algebra is generically not positive: the interpolating categories exist
and are semisimple at generic parameter but are not unitary
\cite{deligne}, and in the Temperley--Lieb instance positivity holds
exactly on the Jones set $\{4\cos^2(\pi/k)\} \cup [4, \infty)$
\cite{jones}.  The existence of a symmetry is thus a positivity
question of the same linear-matrix-inequality type the instrument runs
at every budget: \emph{which} groups can exist is a moment-cone
membership, and the integrality of $N$ is an output.  At generic large
parameter the functional degenerates to the planar (free) functional,
the master field of the large-$N$ limit \cite{voiculescu}; the
certified interpolation between the integer points and the planar
functional is unexplored territory of the same cone.
\end{remark}

\subsection{A structural candidate for the sign problem}
\begin{conjecture}[Dual-pair criterion]\label{conj:signfree}
A fermionic content admits the local positive ground decomposition of
Open Problem~\ref{op:sign} if and only if it is the singlet compression
of a content in the stoquastic class of
Theorem~\ref{thm:factorization}.  The forward half (compressions of
stoquastic content inherit a decomposition on invariant words) is the
tractable half and should be settled first; the converse would make
sign-problem freedom and hidden gauge structure one property.  Either
resolution sharpens Open Problem~\ref{op:sign} from a case list toward
a criterion.
\end{conjecture}

\subsection{Decisive checks, specified in advance}
None of the following has been executed.  (a) \emph{Prepotential
re-pin}: re-derive the certified two-loop brackets of
\S\ref{sec:instance} in prepotential words; the $E_1$ windows and the
metric normalization $\Gamma(x,x) = (4 - x^2)/4$ must reproduce to the
digit, the same moment problem posed in different generators, and any
discrepancy refutes the compression as implemented.  (b) \emph{Casimir
dictionary}: verify $C_2 = (\hat N/2)(\hat N/2 + 1)$ as a
finite-dimensional operator identity per truncation.  (c) \emph{Rank
drop}: the certified kernel of the degree-$\le b$ fusion Gram at
parameter $N = 2$ must coincide with the degree-$\le b$ slice of the
second-fundamental-theorem ideal, the executable form of
Proposition~\ref{prop:groupclosure}.  (d) \emph{Complex-point rank
drop}: check (c) at symmetric pairing and parameter $\delta = 4$; the
certified kernel of the degree-$\le b$ Brauer Gram must coincide with
the degree-$\le b$ slice of the rank-five determinantal ideal of
$O(4,\mathbb C)$, the executable half of
Remark~\ref{rem:complexpoint}, with no Lorentzian input anywhere in
the computation.  Check (a) also bears on
the strong-coupling constant of Appendix~\ref{app:strong}: in
prepotential variables the group-manifold measure is replaced by the
flat Fock measure, so $c_0$ becomes a normal-ordering computation on
polynomial operators---an independent route to the derived value and a
cross-check on the compression itself.

\section{Space from the state}\label{sec:space}

Space is not declared; it must be constructed from $(\A,\w)$.  The known
theorems run from nets of subalgebras to geometry: a half-sided modular
inclusion generates a positive-energy translation group
\cite{wiesbrock,guidolongo};
families satisfying geometric modular action determine spacetimes and
isometry groups \cite{summers}; modular flow of wedge algebras
is geometric in the free field \cite{bisognano}.  What is missing from
the literature, and supplied here as the committed construction, is the
map from $(\A,\w)$ to the net.

\begin{remark}[Motivation for the construction; not a uniqueness
theorem]\label{rem:exhaustion}
The imported theorems consume two structural features: a directed
family of nested subalgebras admitting a common cyclic separating
vector (half-sided inclusions), and a disjointness relation, i.e.\ pairs of
members with commuting or state-independent generation at extensive
content (the spacelike structure of geometric modular action).  Among
the structures the state defines in \S\ref{sec:canonical}, the ball net
of the spectral metric supplies both from outputs alone.  The
alternatives each lack at least one: the Gelfand ideal is a criterion,
not an index; the modular data of the full algebra trivialize on a pure
state; sector gradings are discrete, with no directed region structure;
and the budget filtrations (including multi-parameter ones, which are
genuinely partially ordered) are content-uniform, so distinct levels
are nested and never disjoint, supplying direction without
disjointness (partitioning the generators to manufacture disjointness
re-declares regions, which is the input this section exists to output).
This elimination is a survey of the structures the theory currently
defines, not a theorem that no other index can exist; the construction
below is therefore \emph{committed} rather than derived, and the weight
rests on the decisive tests of Hypothesis~\ref{hyp:gma}, not on this
remark.
\end{remark}

\begin{definition}[The ball net]\label{def:ballnet}
For metric balls $B$ of $(\mathrm{Spec},\, d)$, let $\A(B)$ be the
subalgebra of elements with derivation support in $B$.  $\A(B)$ is
admissible when $\A(B) \cap N_\omega = \{0\}$ (the state separating on
it): the closure structure decides which regions carry geometry.
\end{definition}

\begin{hypothesis}[Geometric modular action of the ball net]
\label{hyp:gma}
The assignment $(\A,\w) \mapsto \{\A(B)\}$ is functorial, and on
extensive content its relative modular data satisfies geometric modular
action, the emergent translations and isometries being those of the
output geometry.  \emph{Scope:} on small bosonic instances the ball net
localizes configuration space and no spatial claim is made; the spatial
claim is made on extensive (many-body, fermionic) content.
\end{hypothesis}

Hypothesis~\ref{hyp:gma} carries four decisive tests in the free chiral
fermion, where every answer is known and no adjustable element exists:
the metric of the vacuum's Dirichlet form must reproduce arc length; the
ball net must reproduce the interval net; the modular flow of a
half-space algebra must be the known dilation flow \cite{bisognano}; and
the translations generated by nested half-space inclusions
\cite{wiesbrock} must be the physical ones.  Failure of any one refutes
the hypothesis.  None of the four has been executed: each requires the
fermionic machinery of Open Problem~\ref{op:car}, which is not yet
built; they are specified in advance so that failure is defined before
the machinery exists.  The division of labor is stated plainly: the
imported theorems are standard, and the entire difficulty lives in
their hypotheses (cyclic separating vectors, half-sided inclusions,
geometric modular action); Hypothesis~\ref{hyp:gma} is exactly the
assertion that the ball net meets those hypotheses on extensive
content.  No instance of emergent geometry, not even a toy, is
exhibited in this paper.

\begin{openproblem}\label{op:car}
Construct the quasi-free (CAR) instance of the Dirichlet-form and
derivation machinery of \S\ref{sec:dirichlet}.  This construction is
prerequisite both to the tests above and to all fermionic content of
\S\ref{sec:hadrons}.  The correspondence of \S\ref{sec:correspondence}
routes this construction and the gauge machinery through one object
(Open Problem~\ref{op:dualpair}).
\end{openproblem}

\subsection{The spacetime group: the real form as modular output}
\label{sec:realform}
The correspondence of \S\ref{sec:correspondence} outputs the spacetime
symmetry only up to complexification: invariant words see the rank and
pairing symmetry of a form, never its signature, so the closure ideal
at parameter $\delta = 4$ fixes $O(4,\mathbb C)$ and cannot distinguish
its real forms (Remark~\ref{rem:complexpoint}).  Real forms of a
complex group are classified by conjugacy classes of antilinear
involutions, and the canonical antilinear data of this theory is the
modular conjugation $J_{\mathcal N}$ of admissible subalgebra--state
pairs (\S\ref{sec:modular}).  The division of the group is then exactly
the division of the instrument: the compact real form is unitarized by
the GNS carrier of the fusion functional, which is positive and tracial
(the reference rung, the Euclidean side); the noncompact generators,
the boosts that distinguish the Lorentz form, are carried by the
modular flow of the ground sector, $\Delta^{it}$ of wedge-type members
of the ball net, which is the content of the Bisognano--Wichmann
theorem on the free field \cite{bisognano} and of geometric modular
action in general \cite{summers}.  On this reading the Lorentz group is
never a declared symmetry: the closure ideal supplies the complex
point, the modular conjugations select the involution, and Wick
rotation is not a contour choice but the KMS analyticity of the modular
read (\S\ref{sec:subsumption}).  Signatures with two causal directions
are excluded internally rather than by declaration: a half-sided
modular inclusion generates positive energy along one causal direction
\cite{wiesbrock}, and no second independent inclusion of modular type
is available, so $O(2,2)$-type forms fail the hypotheses.  The
quaternionic real form $O^*(4)$, which preserves no real symmetric
form, is the homogeneous isometry group of no signature and therefore
never arises as the output of geometric modular action on a metric
ball net; it is excluded at the same step.  The
disconnected components (parity, time reversal) are likewise
antiunitary, $J$-type data, not new inputs.

\begin{definition}[Signature consistency requirement]
\label{def:signature}
On extensive content satisfying Hypothesis~\ref{hyp:gma}, the
complexification of the homogeneous isometry group produced by the
geometric modular action of the ball net must equal the complex group
fixed by the closure ideal of \S\ref{sec:correspondence} at parameter
$\delta$ equal to the read spectral dimension.  Every arrow in the
chain (state to dimension read, dimension to $\delta$, $\delta$ and
ideal to complex point, modular conjugations to real form) is linear in
the state or classical algebra; failure of the equality at certified
resolution refutes the identification, exactly as
Definition~\ref{def:weyl} refutes the geometric one.
\end{definition}

The requirement is conditional on Hypothesis~\ref{hyp:gma} and inherits
its status: nothing here is exhibited, and the comparison is blocked on
Open Problem~\ref{op:car} together with the four tests above, which now
do double duty: the half-space modular flow test is simultaneously the
test that the real form, the signature of spacetime, is an output of
the state rather than an input to the theory.

\section{Extensivity}\label{sec:extensivity}

\subsection{Energy densities}
\begin{theorem}[Certified thermodynamic densities]\label{thm:density}
Let $H_\Lambda = \sum_{\ell \in \Lambda} C_2^{(\ell)} - \lambda
\sum_{p \subset \Lambda} \chi_p$ on boxes $\Lambda$ (interior
plaquettes), $G$ compact, with $c_\chi := \|\chi_p\|$ (for the SU(2)
fundamental character $c_\chi = 2$; in general the dimension of the
plaquette representation).  Then $e(\lambda) = \lim
E_0(\Lambda)/|\Lambda|$ exists, and for every $\Lambda$,
\[
\Big|\, e(\lambda) - E_0(\Lambda)/|\Lambda| \,\Big| \;\le\;
\frac{c_\chi\lambda\, n_\partial(\Lambda)}{|\Lambda|},
\]
with $n_\partial$ the exact count of boundary-crossing plaquettes per
tile.  Certified finite-volume brackets on $E_0(\Lambda)$ therefore yield
certified brackets on $e(\lambda)$ of width $(U - L)/|\Lambda| +
2c_\chi\lambda\, n_\partial/|\Lambda|$.
\end{theorem}
\begin{proof}
Tile $N\Lambda$ by $N$ copies: electric terms partition exactly over
links; magnetic terms split into interior and crossing, $H_{N\Lambda} =
\sum_i H_\Lambda^{(i)} + W_N$, where $W_N$ collects the crossing
plaquettes, of which there are at most $N n_\partial$ (each crossing
plaquette assigned to one of the tiles it meets), so $\|W_N\| \le
c_\chi\lambda N n_\partial$.
Hence $E_0(N\Lambda) \ge N E_0(\Lambda) - c_\chi\lambda N n_\partial$
as an
operator bound, and evaluating $H_{N\Lambda}$ in the product of copy
ground states, $E_0(N\Lambda) \le N E_0(\Lambda) + c_\chi\lambda N
n_\partial$.  For existence of the limit, apply the same interface
estimate to the bipartition of a box into two sub-boxes: it shows
$E_0(\Lambda) + c_\chi\lambda\, n_\partial(\Lambda)$ is subadditive
under box concatenation along each axis, so Fekete's lemma gives
convergence of $E_0(\Lambda_k)/|\Lambda_k|$ along side-doubling
sequences; the two displayed operator bounds then sandwich an
arbitrary box between the tilings it embeds in and those that embed in
it, with per-volume error $c_\chi\lambda\, n_\partial/|\Lambda|$,
which transfers the limit to arbitrary boxes and identifies it as
box-shape independent.  Dividing the two bounds by $N|\Lambda|$ and
letting $N \to \infty$ gives both inequalities with the same constant.
\end{proof}

\subsection{Gaps}\label{sec:gaps}
On the stoquastic class of Theorem~\ref{thm:factorization}, which
includes pure gauge content, $h - E_0 = \sum_x Q_x$ with $Q_x$ local,
positive, and $\Omega \in \ker Q_x$ for all $x$: the ground carrier
automatically presents the structure that finite-size gap criteria
require \cite{knabe,nachtergaele}.  The quantitative content of those
criteria is an overlap (angle) condition between local kernels, which is
a moment read (Gram data) and hence certifiable.  For fermionic
content the precondition is not automatic and may be obstructed (Open
Problem~\ref{op:sign}); the criterion of this paragraph does not reach
such content.  One transport step remains open on the stoquastic class:

\begin{openproblem}[Carrier transport]\label{op:carrier}
The decomposition of Theorem~\ref{thm:factorization} is taken in the
carrier of the ground state of each volume.  Establish certified bounds
transporting the kernel angles of the decomposition on a larger volume
from moment reads on a smaller one together with a certified overlap of
the two ground states on local elements.
\end{openproblem}

\noindent Given Open Problem~\ref{op:carrier}, the martingale argument
of \cite{nachtergaele} yields a certified finite-size gap criterion on
the stoquastic class: certified block gaps and certified kernel angles
imply an explicit infinite-volume gap bound.  The strong-coupling route
is independent of the factorization and is statistics-blind: the
unperturbed part, electric Casimirs per link together with fermion mass
terms, is a commuting, gapped local Hamiltonian with product ground
state, and the magnetic and hopping terms are sums of uniformly bounded
local perturbations, so uniform-in-volume gap persistence below an
explicit threshold follows from the stability theory of gapped
Hamiltonians \cite{bhm} with or without fermions; making the constants
of that theory explicit for the present content is mechanical and
remains to be carried out.  Finite-volume mass shifts in charged sectors
are exponentially small in the volume once the neutral gap and
clustering are certified, in the standard manner \cite{luscher}.
Summarizing the reach: at strong coupling, certified infinite-volume gap
control is available for all content; at moderate coupling it is
available on the stoquastic class conditionally on Open
Problem~\ref{op:carrier}, and for fermionic content it is open at the
structural level (Open Problem~\ref{op:sign}).

\subsection{Order}\label{sec:order}
Symmetry-restricted minimization does not move a determined per-volume
value: any invariant state is admissible for the restricted problem, and
any restricted minimizer averages over the compact symmetry to an
invariant state of equal per-volume value (the value being affine and
translation covariant).  Crystalline or otherwise ordered structure is
therefore never an input; whether the optimizing state orders is read
from its two-point emissions.

\section{Refinement and the continuum}\label{sec:continuum}

The continuum object is defined relationally: two declared contents are
equivalent when all certified reads of $(\A, \w, D_\omega)$ agree; the
continuum theory is the equivalence class.  This section reduces the
existence of that class, for asymptotically free gauge content, to a
finite certification problem, and states exactly what remains.

\subsection{Refinement is exact}
\begin{proposition}\label{prop:projection}
Budget coarsening acts on moment \emph{bodies} by exact linear
projection: $\mathcal S_b$ is the image of $\mathcal S_{b'}$
($b' > b$) under the restriction of moment data.  Content refinement
(subdivision at declared coupling) composes with this projection
through the embedding of the coarser content's algebra in the finer.
No approximation enters the map; all dynamical content of the
continuum question resides in the behavior of certified brackets along
it.
\end{proposition}
\begin{proof}
Restriction maps $\mathcal S_{b'}$ into $\mathcal S_b$ because a
state's degree-$\le b'$ data restrict to its degree-$\le b$ data;
surjectivity is immediate because every point of $\mathcal S_b$ is by
definition the restriction of some admissible state, whose
degree-$\le b'$ data supply a preimage.  For composition with content
refinement: a state of the finer content restricts along the unital
embedding to a state of the coarser, and conversely every state of the
coarser content extends to the finer one (positive Hahn--Banach
extension in the order-unit norm), so the coarse moment body is the
projection of the fine one.
\end{proof}

\begin{remark}[The relaxations do not project exactly]
\label{rem:projectionscope}
The corresponding statement for the outer relaxations $\mathcal M_b$
of \S\ref{sec:cones} is false in general: a level-$b$ feasible moment
vector need not extend to a level-$b'$ feasible one.  Indeed, if it
always did, a diagonal extraction along $b'$ (using the archimedean
moment bounds of the compact generators, as in
Theorem~\ref{thm:totaldichotomy}) would exhibit every level-$b$
feasible point as the restriction of a genuine state, making every
finite level exact and every bracket close at every degree,
contradicting Theorem~\ref{thm:closure}.  The failure of the
relaxations to project exactly \emph{is} the tightening of the
hierarchy; the exactness claimed by
Proposition~\ref{prop:projection} is a property of the bodies, and it
is the bodies along which the trajectory of
Definition~\ref{def:trajectory} is defined.
\end{remark}

\begin{definition}[Trajectory]\label{def:trajectory}
The \emph{trajectory} is the locus in (content, coupling) along which
per-step contraction certificates exist for all emitted relational
reads: certified nesting
\[
\mathrm{bracket}(n{+}1) \subseteq \mathrm{bracket}(n),
\]
geometric in the directions transverse to the
locus, with the declared logarithmic flow along it (the marginal
direction of an asymptotically free coupling).  The step constraint set
at stage $n$ is the certified bracket set emitted at stage $n{-}1$
together with the content's region constraints; no window parameters
exist.
\end{definition}

\begin{theorem}[Tower summation]\label{thm:tower}
If per-step certified control holds along the trajectory with
$\delta_n \le C\rho^n \operatorname{polylog}(n)$, $\rho < 1$, then the
certified brackets on every emitted read converge to a limit interval
with explicit enclosure $r_N + [-T_N, T_N]$, $T_N = \sum_{n \ge N}
\delta_n$.
\end{theorem}
\begin{proof}
The intervals form a Cauchy sequence in Hausdorff distance with exactly
summable tail; the enclosure is the triangle inequality summed from $N$.
\end{proof}

\subsection{Two dimensions are exact}
\begin{theorem}\label{thm:2d}
In two dimensions with heat-kernel plaquette weights
\[
K_t(U) = \sum_j d_j e^{-t C_2(j)} \chi_j(U),
\]
subdivision of a plaquette of area $t$
into areas $t_1 + t_2 = t$ leaves every gauge-invariant read exactly
invariant.
\end{theorem}
\begin{proof}
Character orthogonality gives $\int \chi_j(UV^\dagger)\chi_k(VW^\dagger)
\,dV = \delta_{jk}\,\chi_j(UW^\dagger)/d_j$, hence $K_{t_1} * K_{t_2} =
K_{t_1+t_2}$: the subdivided theory equals the original on the invariant
algebra, and invariant reads depend on enclosed areas alone.
\end{proof}

\noindent Theorem~\ref{thm:2d} is the heat-kernel form of the exact
solvability of two-dimensional lattice gauge theory \cite{migdal},
transported for the dependency graph; it is the degenerate case of
Definition~\ref{def:trajectory} (contraction ratio identically zero).
At the opposite end of solvability, hierarchical models and free fields
constitute the calibration class: certified contraction rates computed
there are required to equal the known renormalization-group eigenvalues
of those models; rigorous computer-assisted renormalization is an
established genre \cite{lanford,koch}, of which the present certificates
are the conic instance.  A construction that produced rates in
disagreement with the known spectra would thereby be refuted as a
construction, whatever its internal consistency.

\subsection{Slack, dichotomy, and the anchor}
\begin{definition}[Slack]\label{def:slack}
For a state $\omega$, $\sigma(\omega)$ is the supremum, over certificates
of the per-step inequality family at fixed degree, of the minimal margin.
As a supremum of affine weak-$*$ continuous functionals, $\sigma$ is
weak-$*$ lower semicontinuous and nonnegative on feasible states.
\end{definition}

\begin{lemma}[Uniformity on compacts]\label{lem:uniformity}
If every state in a weak-$*$ compact set $K$ admits a certificate of
degree $\le b$ with slack $\ge s > 0$, then finitely many certificates
cover $K$ and a uniform budget certifies the step on all of $K$.
\end{lemma}
\begin{proof}
A fixed certificate's margin is affine and weak-$*$ continuous, so slack
$\ge s$ at a point gives margin $> s/2$ on a neighborhood; compactness
(Banach--Alaoglu) extracts a finite subcover.
\end{proof}

\begin{lemma}[Neighborhood persistence]\label{lem:persistence}
If $\sigma(\omega^*) = s^* > 0$ then $\sigma \ge s^*/2$ on a weak-$*$
open neighborhood of $\omega^*$.
\end{lemma}
\begin{proof}
The certificate achieving $s^*$ has affine continuous margin; its
superlevel set at $s^*/2$ is open and lower-bounds $\sigma$.
\end{proof}

\begin{theorem}[Dichotomy]\label{thm:dichotomy}
Along any refinement orbit $\{\omega_n\}$, either $\inf_n
\sigma(\omega_n) > 0$, or the orbit has a weak-$*$ accumulation point
$\omega^*$ with $\sigma(\omega^*) = 0$ exactly.
\end{theorem}
\begin{proof}
If $\liminf \sigma(\omega_n) = 0$, extract a subsequence with $\sigma \to
0$; it accumulates by compactness of the state space; lower
semicontinuity gives $\sigma(\omega^*) \le 0$, and nonnegativity closes.
\end{proof}

\noindent A zero-slack state is one at which a transverse direction
ceases to contract: an unanticipated criticality.  The classical failure
mode of rigorous renormalization, uncontrolled flow away from the
governing fixed point, is here \emph{derived} as the unique failure
mode.

The step map factorizes as (variational update) $\circ$ (restriction
channel).  The channel half is monotone for relative entropy
\cite{uhlmann}, and by the variational formula of \cite{kosaki} relative
entropy is itself a supremum of affine functionals of the state, hence an
object of the same class as $\sigma$; the variational half-step is where
all difficulty resides.

For asymptotically free content the ultraviolet tower flows toward the
free fixed point, where the theory is exactly solvable and the slack of
the per-step family, with the coupling retained as a bounded parameter,
is a finite computation on free content.  If that slack is positive,
Lemma~\ref{lem:persistence} yields an open neighborhood $N$ of the free
profile on which a single \emph{parametric} certificate (multipliers
polynomial in $g$ over $[0, g^*]$, in the standard manner of
\cite{putinar,lasserre}) can certify every step, including the
forward-invariance inequality keeping the orbit in $N$; near the free
anchor, effective degree bounds for such certificates \cite{niesch}
render the search decidable with explicit budget.  If the free-field
slack is not positive at reasonable degree, the strategy of this section
fails at its first computation; the question is settled either way by
the computation, which has not yet been carried out.

\subsection{The reduction}
\begin{theorem}[Reduction]\label{thm:reduction}
Suppose (i) a parametric certificate with positive slack and certified
forward invariance exists on a neighborhood $N$ of the free profile,
uniformly for $g \in [0, g^*]$; and (ii) a finite chain of per-step
contraction certificates connects the declared content to $N$.  Then the
certified brackets on every emitted relational read converge along the
tower with explicit enclosure and uniform budget: the continuum brackets
exist, and their existence is verified by finitely many exact-arithmetic
identities.
\end{theorem}
\begin{proof}
The finitely many chain steps carry exhibited certificates; steps beyond
the chain are covered by the parametric certificate, whose invariance
clause keeps the orbit in $N$ and whose contraction clause supplies
$\delta_n$ as in Theorem~\ref{thm:tower}; that theorem closes the
enclosure.  The budget is the maximum over a finite set.
\end{proof}

\begin{remark}[What the reduction changes]\label{rem:teeth}
Theorem~\ref{thm:reduction} changes the \emph{type} of the continuum
problem, from a limit statement in analysis to a finite certification,
conditional on its hypotheses; it does not by itself reduce the
difficulty, which is relocated intact into hypothesis (ii), the chain
of Conjecture~\ref{conj:chain}, and into hypothesis (i), whose
free-anchor slack computation has not been carried out.  The claim that
the certificate language has teeth is earned only by executions; the
unexecuted decisive computations are named in \S\ref{sec:conclusion}.
\end{remark}

\subsection{Both branches have witnesses; data cannot decide}
\begin{lemma}[Duality]\label{lem:duality}
On content with compact generators, the per-step feasibility problems at
degree $\le b$ (constraint slices being certified brackets of positive
width together with the region constraints) admit infeasibility
witnesses: if no certificate exists, there is a normalized functional
$\Lambda_b$, nonnegative on the truncated quadratic module.
\end{lemma}
\begin{proof}[Proof sketch]
The slices have nonempty interior, so the truncated quadratic modules
are closed \cite{powers,marshall}; conic separation yields the witness.
For the interior point: the Haar state alone is not feasible for the
slices (its moments generally violate the emitted brackets; in the
worked instance $\langle x{+}y{+}z\rangle_{\mathrm{Haar}} = 0$ against
a certified lower bound above $4$), but a convex combination
$(1-\varepsilon)\,\omega + \varepsilon\,\omega_{\mathrm{Haar}}$ of a
slice-feasible state with the Haar state is: faithfulness of the Haar
state on compact content makes its moment matrices strictly positive
definite, so the mixture is strictly interior to the positivity
constraints, and the positive width of the bracket slices keeps the
mixture inside them for $\varepsilon$ small.  The Slater alternative
then applies.
\end{proof}

\begin{theorem}[Localization; total dichotomy]\label{thm:totaldichotomy}
Fixed-degree certifiability is open in the coupling (a certificate with
slack absorbs $O(\delta g)$ coefficient perturbations into its margin),
so $F_b = \{g \in [0, g_c] : \text{step certificate of degree } \le b\}$
is open and increasing in $b$.  Consequently either $\bigcup_b F_b =
[0, g_c]$ and compactness extracts a finite chain
(Theorem~\ref{thm:reduction}(ii) holds); or the residue $[0,g_c]
\setminus \bigcup_b F_b$ is nonempty and compact, and at each of its
points the witnesses $\Lambda_b$ of Lemma~\ref{lem:duality} exist for
every $b$ and accumulate weak-$*$ on genuine states of exactly zero
slack.  Either the chain exists, or a zero-slack obstruction state
exists at a localized coupling: both branches of the question have
witnesses.
\end{theorem}
\begin{proof}
Openness: fix a certificate with margin $s$ at $g_0$; the inequality
coefficients are polynomial in $g$, so for $|g - g_0|$ small the same
certificate certifies the perturbed inequality with margin
$s - O(|g - g_0|)$ by the residual-norm clause of
Proposition~\ref{prop:weak}; hence $F_b$ is open, and increasing in $b$
because certificates persist under degree enlargement.  If the $F_b$
cover $[0, g_c]$, compactness extracts a finite subcover, hence a
finite chain.  Otherwise fix $g^\dagger$ in the compact residue; for
every $b$ a witness $\Lambda_b$ exists at $g^\dagger$
(Lemma~\ref{lem:duality}).  Each $\Lambda_b$ is normalized and positive
on squares of degree $\le b$, and Cauchy--Schwarz on the truncated
module together with the exact norm bounds of the compact generators
gives $|\Lambda_b(w)| \le 2^{\deg w}$ for every word of degree
$\le b$.  For each fixed degree $d$, the restrictions of
$\{\Lambda_b\}_{b \ge d}$ to the degree-$d$ word space therefore lie in
a closed bounded, hence compact, set; a diagonal extraction along $b$
yields a functional $\Lambda$ on all words, normalized, positive on
squares of every degree, with $|\Lambda(w)| \le 2^{\deg w}$.  That
bound makes $\Lambda$ norm-continuous on the dense polynomial
subalgebra, so it extends to a state on the C$^*$ closure.  Finally,
each $\Lambda_b$ assigns the step inequality a nonpositive value while
satisfying every slice constraint (that is the witness property); the
finitely many margins are jointly continuous in $(\Lambda, g)$, so the
limit state assigns the step inequality a nonpositive value while lying
in the constraint slice: a state of exactly zero slack at
$g^\dagger$, since positive slack at the limit would, by
Lemma~\ref{lem:persistence}, certify a weak-$*$ neighborhood meeting
the $\Lambda_b$ and contradict their witness property.
\end{proof}

\begin{theorem}[No decision by data]\label{thm:nogo}
The witnesses $\Lambda_b$ satisfy every slice constraint by
construction (hence every certified bracket the instrument has emitted
or will emit on the declared reads, up to their degree), and so does any
limiting obstruction state.  Consequently no finite accumulation of
certified brackets on the physical state entails the nonexistence of the
obstruction: the branches of Theorem~\ref{thm:totaldichotomy} cannot be
separated by bracket data alone.  They can be separated only by
exhibition of the chain, or by a structural theorem excluding the
obstruction.
\end{theorem}

\begin{remark}[What failure looks like]
By Lemma~\ref{lem:duality}, failure at fixed degree is itself decidable
and carries a witness; a sustained failure across degrees produces a
sequence of witness functionals converging on the obstruction's own
moment profile, which can be symmetrized and interrogated by every
declared read.  Unsuccessful search at bounded degree is therefore
informative; unbounded non-existence is not finitely refutable, which is
the precise character of the remaining problem.
\end{remark}

\subsection{The central conjecture and candidate mechanisms}
\begin{conjecture}[Existence of the chain]\label{conj:chain}
For asymptotically free gauge content in $3{+}1$ dimensions, the finite
chain of Theorem~\ref{thm:reduction}(ii) exists.  Equivalently
(Theorem~\ref{thm:totaldichotomy}): no zero-slack obstruction state lies
on the segment between the declared content and the free anchor.  This
is nonperturbative asymptotic freedom in certificate form; upon
exhibition of the chain, the continuum brackets of the theory exist at
the strength of finitely many verified identities.
\end{conjecture}

Three mechanisms, each with a decisive computation in the calibration
class, would produce or organize the chain.  \emph{(a) Curvature
condition.}  The canonical per-step certificate is the running state's
own geometry: transverse contraction as a uniform spectral gap of the
state Dirichlet forms, certified by a curvature-dimension inequality
$\Gamma_2 \ge \kappa\,\Gamma$ \cite{bakry} in the exact $\Gamma$-calculus
of \S\ref{sec:metricdim}; whether the per-step certificates of the
solvable controls can be taken in this form is a computation.  On
fermionic content this mechanism inherits the scope boundary of
Remark~\ref{rem:scope} and Open Problem~\ref{op:sign}.  \emph{(b)
Monotone.}  A certified functional of the state, in the affine-supremum
class of Definition~\ref{def:slack}, strictly decreasing along the
refinement flow with quantitative strictness away from the free fixed
point, would exclude the obstruction branch outright, and with the free
anchor computation would imply Conjecture~\ref{conj:chain} by
Theorem~\ref{thm:totaldichotomy}; the entropic proofs of
renormalization-flow monotonicity \cite{casini} indicate the class in
which to look, and the requirement here, a monotone for one flow on one
content class, is strictly weaker than a universal theorem.  \emph{(c)
Parametric collapse.}  If the chain's certificates can be taken with
coefficients polynomial in $g$ over the whole segment, the chain is a
single parametric certificate; whether the per-step certificates of the
controls vary polynomially along the flow is a measurement.

\section{Hadronic content}\label{sec:hadrons}

\subsection{Fermions at declared content}
Quark content enters as CAR sectors coupled to the gauge content, the
Gauss law imposed by exact quotient on the invariant algebra (at declared
lattice content the quotient is exact and requires no section, so the
obstruction that blocks continuum nonabelian gauge fixing does not arise
at this level).  Lattice chirality is constrained: no discretization
carries exact chiral symmetry, doubler freedom, and locality
simultaneously \cite{nielsen}; the class defined by the Ginsparg--Wilson
relation $\{\gamma_5, D\} = a D\gamma_5 D$ \cite{ginsparg,luscherchiral}
carries the exact algebraic form of chirality and is adopted as content.
Residual parameters within the class are declared, and every emission is
required to be stable across the declared parameter window at its stated
resolution; parameter dependence at resolution refutes the emission, not
the framework.

\subsection{Spectral hadrons}
Baryon and flavor numbers are canonical gradings of the content
(\S\ref{sec:metricdim}); the mass of the lightest baryon is the
ground-sector energy difference between grading sectors,
\[
m_B \;=\; E_0(B{=}1) - E_0(B{=}0),
\]
a spectral quantity requiring no interpolating operators, no plateau
analysis, and no statistical estimator: both terms are certified
brackets, and the difference inherits interval arithmetic.  Emissions
are dimensionless ratios (e.g.\ $m_B/m_M$ against the lightest meson);
each content parameter (quark-mass coupling) is fixed by exactly one
declared calibration ratio, disjoint from all comparison targets, and
electromagnetic and isospin-breaking effects are outside the stated
resolution tier until declared otherwise.

\subsection{What a hadronic emission requires}
A certified $m_B/m_M$ at declared lattice content and coupling requires:
the fermionic instance of the canonical machinery (Open
Problem~\ref{op:car}); certified brackets in two grading sectors on
volumes governed by \S\ref{sec:extensivity}, where the reach is exactly
as summarized in \S\ref{sec:gaps}: Theorem~\ref{thm:density} is
statistics-blind (fermionic hopping terms crossing a tile boundary are
norm-bounded per link and mass terms partition, so the proof extends
verbatim), the strong-coupling gap control is statistics-blind, and the
moderate-coupling gap criterion does \emph{not} reach fermionic content
(Open Problem~\ref{op:sign}); and, for the comparison to carry weight,
the strong-coupling value of the ratio derived independently in advance,
with the resolution rule of \S\ref{sec:resolving} in force.  The
strong-coupling lattice baryon program is therefore reachable by the
statistics-blind route; away from strong coupling, certified fermionic
gap control is open at the structural level.  A certified
\emph{continuum} $m_B/m_M$, the proton-to-meson ratio of the physical
theory, requires in addition exactly the content of
\S\ref{sec:continuum}: it exists upon Conjecture~\ref{conj:chain}
together with the free-anchor computation, by
Theorem~\ref{thm:reduction}, and by Theorem~\ref{thm:nogo} no
accumulation of lattice data can substitute for that exhibition.  The
open items between the present state of the framework and a physical
proton mass are therefore precisely: Open Problems \ref{op:car},
\ref{op:carrier}, and \ref{op:sign}, the explicit constants of
\S\ref{sec:gaps}, the free-anchor computation, and
Conjecture~\ref{conj:chain}.  These are not of one kind: all but one are
completed by work of a known character, while Open
Problem~\ref{op:sign}, the sign problem in carrier form, may be a
structural obstruction rather than an unfinished task, and any account
of the path that omits it, or lists it as routine, misstates the
theory's position.  Nothing else stands in the path, and each item is
stated where its mathematics lives.

\section{Electrodynamic matter}
\label{sec:subsumption}

This section declares the content of electrodynamic matter, with
gravitation carried as a spin-2 boson sector, and develops it inside the
framework.  Historically this instance was constructed first; the general
theory of \S\S\ref{sec:kinematics}--\ref{sec:continuum} is its closure,
and everything below is that theory evaluated at declared content.

\subsection{Content}\label{sec:qedcontent}
The arena is $\mathbb R^3$, entering the theory exclusively through the
coefficient data of the local dynamics (kernels, dispersions, form
factors); this is the entire residual background.  A system is declared
by its composition: the species present, a closed subgroup $\mathsf H
\subseteq \mathbb R^3$ (the \emph{content group}; every closed subgroup
is $V \oplus \Lambda$ with $V$ a subspace and $\Lambda$ a lattice in a
complement) under which the matter distribution is homogeneous, and the
species densities on $\mathbb R^3/\mathsf H$: bulk matter is $\mathsf H =
\mathbb R^3$, a molecule $\mathsf H = \{0\}$.  The content group declares
what system is studied, never how its matter organizes; no lattice,
cell, or space group is among the data (\S\ref{sec:order}).

Species are second-quantized fields.  Per charged-lepton flavor, the CAR
algebra over
\[
L^2(\mathbb R^3) \otimes \mathbb C^4
\]
graded by electric
charge: particles and antiparticles carried together, particle number
not conserved, flavor charge conserved, and \emph{no positive-energy
subspace selected at any point for any purpose}
(\S\ref{sec:sea-instance}).  Per declared isotope, a CAR or Weyl-form CCR
field over $L^2(\mathbb R^3) \otimes \mathbb C^{2s+1}$ with mass, charge,
spin, and charge form factor as Hamiltonian data of the standing of the
mass.  The photon is the transverse Weyl algebra at bandwidth (two
polarizations).  The graviton is the transverse-traceless spin-2 Weyl
sector, coupled to the matter stress words exactly as the photon couples
to the charge current, with $\kappa = \sqrt{16\pi G}$ a measured datum of
the standing of a charge; no separate gravitational theory is carried.
Statistics are algebra identities; no symmetrization machinery exists.
Admissible states are positive, normalized, $\mathsf H$-invariant,
invariant under each flavor's charge $U(1)$ at the declared densities,
neutral in total, with finite boson second moments; per-covolume
quantities live in the density quotient of local elements modulo
translates.  The Hamiltonian density assembles the Dirac coefficient
matrices, the isotope dispersions, the instantaneous Coulomb interaction
of the \emph{total} charge density (produced by the quotient below), the
transverse current coupling, the free field energies at budget, and the
transverse-traceless stress coupling.

\begin{proposition}[Boundedness below at budget]\label{prop:qedbounded}
At every stated budget, $\hat h \ge -C(\mathrm{budget})\,\mathbf 1$ as a
form on the admissible class.
\end{proposition}
\begin{proof}
Fermionic smeared words with bounded coefficients are bounded; kinetic
and boson number forms are nonnegative; the reciprocal Coulomb part is a
positive superposition of Hermitian squares of smeared total-charge
words; the singular slice of the point-charge specialization is
form-dominated by a declared fraction of the lepton kinetic bilinear
(Hardy), itself bounded at bandwidth; and boson-linear couplings obey,
mode by mode, $\pm(Ba^\dagger + B^\dagger a) \le \varepsilon\omega\,
a^\dagger a + (\varepsilon\omega)^{-1} B^\dagger B$ with $B$ bounded,
the first term absorbed into a fraction of the free energy over the
budgeted band.
\end{proof}

\begin{proposition}[Two-sided transfer of kinetic truncation]
\label{prop:kintrunc}
If $\hat h^{(t)}$ carries, for any subset of massive species, kinetic
dispersions truncated from below ($0 \le t(k) \le k^2/2m$ pointwise), then
$\hat h^{(t)} \le \hat h$ as forms: every certified lower bound for the
truncated theory transfers, while every exhibited state evaluates the full
dispersions exactly.  Kinetic truncation costs width only, on both sides,
never soundness.
\end{proposition}

\begin{proposition}[The photon axis is variational from above]
\label{prop:photonaxis}
$e_0$ is monotone nonincreasing in the photon budget: a budgeted admissible
state extended by the vacuum on the new band evaluates the enlarged
Hamiltonian to the same energy density.  Enlarging the budget enlarges the
coupling, which is indefinite, so no form inequality exists; lower bounds
are per-budget, and the behavior of anchored relational emissions along the
axis is the measured trajectory of \S\ref{sec:uvaxis}.
\end{proposition}

\subsection{The exact abelian quotient}
The electromagnetic phase space is reduced, never gauge fixed: the
abelian Gauss constraint is solved exactly on $\mathbb R^3$, the
longitudinal electric field is determined by the charge density, its
elimination produces the instantaneous Coulomb interaction, and the
remaining local degrees of freedom are the two transverse polarizations.
This is a quotient with nothing selected; at $\mathsf H = \mathbb R^3$ or
$\{0\}$ there are no cycles, hence no holonomies and no flux labels.  It is
the $G = U(1)$ case of the exact quotient of \S\ref{sec:content}.

\begin{remark}[Torus holonomies at periodic content]\label{rem:holonomy}
At a content group $\mathsf H = V \oplus \Lambda$ with a nontrivial lattice
$\Lambda$ of rank $d \ge 1$ the reduced configuration space carries, besides
the local transverse modes, the finitely many global harmonic modes of the
periodic geometry: the holonomies of the connection around the $d$ cycles
with their conjugate momenta, and at $d \ge 2$ integer magnetic flux labels
through the $2$-cycles.  Both dispositions are exact and cost nothing.  The
flux labels are superselected sector data, declared with the particle
content; the holonomy pairs are a fixed finite set of global Weyl generators
whose contribution to every per-covolume density vanishes in the quotient,
so no emission sees them until a flux-threading emission is wanted, when they
are adjoined explicitly.  The blanket absence above holds only at the two
endpoint content groups.
\end{remark}

\begin{remark}[Constraints without sections]\label{rem:sections}
Where a constraint's solution space is not a subalgebra (the nonabelian
gauge orbit, which admits no global section), this framework forms no
projection at all: content is declared on the invariant algebra
directly, as in \S\S\ref{sec:instance},~\ref{sec:hadrons}, and the
constrained non-subalgebra projection never arises as an object of the
theory.  The continuum question for the nonabelian sectors is then
exactly Conjecture~\ref{conj:chain}.
\end{remark}

\subsection{The physical inner product and its reads}
\begin{definition}[Single emission of the instance]\label{def:pip}
The \emph{physical inner product} of the content is $\langle A, B
\rangle = \w(A^\dagger B)$ on the word spaces: the GNS form of
\S\ref{sec:canonical} and the moment functional of \S\ref{sec:cones},
carried projectively (the state normalization is gauge).  Every
certified number of the instance is a two-sided bracket on a matrix
element of this form over the ground sector, taken directly or through a
conditional expectation onto a translation-covariant subalgebra (which
preserves both positivity and the per-covolume average, and is the only
conditioning class admitted).
\end{definition}

The energy axis carries an affine gauge, the additive zero point and the
multiplicative clock rate, whose invariants are exactly the ratios of
sector differences.  The absolute energy is evaluation-internal; every
emitted physical quantity is such a ratio, with the reference transition
a measured datum that doubles as a calibration anchor.  The remaining
reads: sector differences at distinct declared compositions; the
structure factor $\w(\hat\rho(k)^\dagger\hat\rho(k))$, whose Bragg atoms
are crystalline order read and never imposed (\S\ref{sec:order}); the
frame, read from the vanishing of the momentum-density one-point; modular
time, the admissible-subalgebra modular structure of \S\ref{sec:modular},
exactly evaluable for quasi-free and Gaussian restrictions under a
full-rank condition on the restricted covariance that fails shut on
degeneracy; and the response family, the jet of the certified energy
tube under source couplings $\hat h \mapsto \hat h + \sum_\alpha
\theta_\alpha O_\alpha$: the gradient is the one-point read (order
parameters and frame; Hellmann--Feynman at stationarity, envelope
one-sided derivatives in general, degenerate optima failing shut), the
Hessian the static susceptibilities, and time-dependent sources the
complex-time two-point whose analyticity is the modular read.

\subsection{The polarization is an output}\label{sec:sea-instance}
The first-quantized Dirac--Coulomb many-body problem is unbounded below
\cite{brownravenhall}.  Confronted with it, a two-sided instrument fails
shut: no certificate exists, which is the true statement, and
variational collapse (a disease of one-sided methods) cannot occur.
On the charge-graded field algebra no positive-energy subspace is ever
selected: the polarization, the projector that no-pair models fix by
hand in representation-dependent ways, is here the level-one moment
datum of the state, an output on the same footing as the nuclear
covariance (phonons), the photon covariance and displacement (the
classical field), and the Bragg content (the crystal).  These are the
level-one reads of the reference rung (Proposition~\ref{prop:reference})
at the field group, the quasi-free$\otimes$Gaussian data the singlet
fusion determines, with the modular generator (the flow) the fifth.  The
stability
literature corroborates the direction: the no-pair model is unstable
with the bare-Dirac electron subspace \cite{lss} and stable with the
field-dressed one \cite{liebloss}; the physically correct polarization
was never a prior selection but a property of the interacting state,
which is where the physical inner product places it.  Ground states of
the nonrelativistic quantized-field models exist at all coupling
\cite{gll}, supplying the exhibited-state side of the brackets.

\subsection{Specializations and calibration}
Four theories, each a coefficient substitution with a declared budget
restriction in the one above, never an approximation within the finer
theory: full quantum electrodynamics at cutoff; the Pauli--Fierz theory
(pair sectors frozen, nonrelativistic kinetic coefficients); the Coulomb
theory (photon budget empty); the clamped theory (isotopes replaced by
classical point sources at parametric positions).  Certified intervals
of adjacent theories at shared budgets bracket, telescopically, the
nuclear-motion content, the radiative and retardation content, and the
relativistic and pair content; no endpoint is estimated, and in
relational form the brackets remain certified and dimensionless.
Isotope shifts decompose into certified field and motional parts, whose
reversal reads nuclear charge radii.  Renormalization is budget-running
data: one measured anchor per counterterm, no anchor anywhere else, each
bare parameter an interval-valued datum propagated by interval
arithmetic; calibration costs width, never soundness.

\begin{proposition}[Certified approximation brackets]\label{prop:brackets}
Certified intervals of adjacent theories at shared budgets bracket, no
endpoint estimated: the nuclear-motion content in $[\,L_{\mathrm{Coul}} -
U_{\mathrm{Cl}},\, U_{\mathrm{Coul}} - L_{\mathrm{Cl}}\,]$, the radiative and
retardation content in $[\,L_{\mathrm{PF}} - U_{\mathrm{Coul}},\,
U_{\mathrm{PF}} - L_{\mathrm{Coul}}\,]$, and the relativistic and pair
content in $[\,L_{\mathrm{QED}} - U_{\mathrm{PF}},\, U_{\mathrm{QED}} -
L_{\mathrm{PF}}\,]$.  The brackets compose telescopically, and in relational
form they divide by the reference charge and stay certified and
dimensionless.  Isotope structure sharpens along the ladder: clamped
energies carry the form factors but not the masses, so their isotope
differences are certified \emph{field} shifts; joint emissions carry both,
so certified isotope shifts decompose into field and motional parts, the
reversal reading nuclear charge radii, with the radiative bracket the
muonic-hydrogen read.
\end{proposition}

\subsection{Convergence for this content}
\begin{theorem}[Finite range]\label{thm:finiterange}
For $\mathsf H$-covariant Hamiltonian densities of finite range
satisfying Proposition~\ref{prop:qedbounded}, the certificate hierarchy
converges: $\sup_\beta L^* = e_0$.
\end{theorem}
\begin{proof}[Proof outline]
Tile $\mathsf H$ by overlapping weighted boxes; on each box the shifted
cluster Hamiltonian, tilted by density multipliers, is a positive
element, hence a Hermitian square whose factor is approximable by
cluster words with certified remainder; cluster grand energies converge
per covolume, and Legendre duality with convexity of the energy density
in the densities, known for the Coulomb class \cite{lieblebowitz},
closes the argument.  The functional-analytic control of the square-root
approximation on the infinite-dimensional bosonic box sectors is the one
step not carried out here (Open Problem~\ref{op:qedconv}).
\end{proof}

For the instantaneous interaction, the identity of Fefferman and
de la Llave \cite{fdll},
\[
\frac{1}{|x-y|} \;=\; \frac1\pi \int_0^\infty \frac{dr}{r^5}
\int_{\mathbb R^3} dz\; \mathbf 1_{B(z,r)}(x)\,\mathbf 1_{B(z,r)}(y),
\]
decomposes the total charge--charge interaction exactly into a
finite-range singular part ($r < r_0$, absorbed into the cluster energies),
a bounded one-body remainder, and for $r \ge r_0$ a part manifestly a
positive superposition of Hermitian squares of smeared \emph{total}
ball-charge words $\hat Q_B$, species-blind by construction, every species a
field density.  With a spectral guard $\hat G(k) \ge g(r_0)$ on the low
modes $k < 1/r_0$ (with $g = c_3 r_0^2,\, c_2 r_0,\, c_1 \ln r_0$ by
noncompact dimension $3, 2, 1$), guard manipulations stay in the
certificate class.

\begin{lemma}[Mass bound]\label{lem:mass}
For every admissible $\omega$ within unit energy of the remainder ground,
$\int_{k<1/r_0}\hat G\,d\hat C_\omega \le C_0$, hence $\int_{k<1/r_0}
d\hat C_\omega \le C_0/g$: the low-mode charge fluctuation the guard admits
is $O(1/g)$.
\end{lemma}

\begin{lemma}[Phase-space integrability; kernel--measure]\label{lem:phase}
With $W(k)$ the long-wavelength kernel, $W(k)\,dk$ is bounded on the guarded
ball up to a logarithm at noncompact dimension one.  The $k=0$ mode is the
exactly conserved total charge: at $\mathsf H = \{0\}$ sector sharpness
kills its atom outright, and at noncompact $\mathsf H$ the atom is removed
by a conditioning step at its degenerate-best instance (identically
vanishing f-sum), carrying the status of Theorem~\ref{thm:cost}.
\end{lemma}

\begin{theorem}[Cost Theorem]\label{thm:cost}
There are constants $K_1, K_2$ from budget data alone (local operator norms,
the reference f-sum $\sum_s n_s q_s^2/m_s$ over charged species, the
stripped response constant $A$) such that for every admissible $\omega$
within unit energy of the remainder ground there is an admissible
$\tilde\omega$ with remainder expectation at most $K_1 g^{-1}$ larger and
deleted-part expectation at most $K_2 g^{-1}$.  Hence $\Delta_{\mathrm{reopt}}
(r_0) \le K g(r_0)^{-1} \to 0$, and with Theorem~\ref{thm:finiterange} on
the finite-range remainder, $\sup_\beta L^* = e_0$ at every content group,
gapped or metallic, ordered or not.
\end{theorem}
\begin{proof}[Proof outline]
Every object is a functional of the \emph{total} charge: dyadic shells with
smeared total-charge quadratures and reference conjugates, one linear cost
law
\[
\Delta E_j \le \min(\text{squeeze}, \text{condition}) + A(S-B),
\]
iterated squeezing with re-derived mass bounds (Lemma~\ref{lem:mass}),
median-halving conditioning where the shell f-sum is depleted, outer-to-inner
sweeps, $B \sim r_0^{-2}$.  The enumerated items are Open
Problem~\ref{op:qedconv}, and each mechanism is pinned in a closed-form
solvable model in Appendix~\ref{app:battery}.
\end{proof}

\emph{Content geometries.}  At $\mathsf H = \{0\}$ flavor-charge sharpness
strengthens the certificate class: sharp states kill the full one-sided
ideal $(\hat Q_\ell - q_\ell)\A + \A(\hat Q_\ell - q_\ell)$, since
$(\hat Q_\ell - q_\ell)\Omega = 0$ in the carrier.  For noncompact
$\mathsf H$ the constraints are density expectations, sharp nowhere, and
only the scalar multipliers are killed; \emph{admitting the ideal there is
unsound}.  This is the one point at which content geometries differ in the
certificate class.

\begin{proposition}[Convergence at noncompact dimension $1,2$]\label{prop:lowd}
$\sup_\beta L^* = e_0$ at $\mathsf H = \mathbb R, \mathbb R^2$ by the proof
of Theorem~\ref{thm:cost} with the substitutions: tiling and thermodynamic
limit over $\mathsf H$; the Fefferman--de la Llave decomposition unchanged
(it decomposes the interaction in $\mathbb R^3$, independent of the content
group); the diverging guard saturations $g = c_2 r_0,\, c_1 \ln r_0$; and
for polar slabs the slab form of the kernel, the transverse direction an
integral, the dipole term finite and exact.  Rates degrade with decreasing
dimension.
\end{proposition}

\begin{openproblem}\label{op:qedconv}
The enumerated items of Theorems~\ref{thm:finiterange} and~\ref{thm:cost}:
the bosonic-sector square-root approximation; the many-body second-order
bookkeeping and explicit quasi-locality constants of the cost law; the
$k=0$ conditioning step at noncompact content groups
(Lemma~\ref{lem:phase}); the carry-over to the boson-linear (transverse)
couplings, per frequency shell with the photon frequency the guard's
analogue; and the low-dimension constants and slab step of
Proposition~\ref{prop:lowd}.  Each is functional-analytic; none affects the
certified per-budget soundness of any emission.
\end{openproblem}

\begin{theorem}[Clamped exactness at $\mathsf H = \{0\}$]
\label{thm:clamped}
In the clamped theory with $m$ budgeted electronic modes, the bracket
closes at word degree $2m$: $L^* = e_0 = U^*$.
\end{theorem}
\begin{proof}
Monomials of degree $\le 2m$ in the mode operators are a linear basis
of the CAR algebra (which is $4^m$-dimensional), so the degree-$2m$
word space is the whole algebra: the moment data at this budget are
the functional itself, and Gram positivity over a spanning word set is
positivity on every Hermitian square, i.e.\ the moment cone at this
degree \emph{is} the state set.  With $P_{N'}$ the number-sector
projectors (themselves words of degree $\le 2m$), $\hat H -
\sum_{N'} E_0(N')\,P_{N'} \succeq 0$ is a sum of Hermitian squares of
algebra elements, hence a certificate at this degree, giving the lower
value $E_0(N)$ on the declared sector; the exhibited side attains it
at full richness.
\end{proof}

\begin{remark}[The half-degree relaxation]\label{rem:clampeddegree}
Products of two degree-$\le m$ words already span the algebra, so the
level-$m$ relaxation determines every moment of the functional.
Whether its positivity constraints alone are already exact (closure at
word degree $m$ rather than $2m$) is a positivity-transfer question:
determination of moments does not by itself imply exactness of a
relaxation, and positivity on the squares of a spanning-product
subspace is strictly weaker than full positivity in general.  It is
verified where used and not claimed here.
\end{remark}

\subsection{Gravitational reads}\label{sec:gravreads}
Every measured gravitational quantity is a read of the physical inner
product; none is an addition to the machinery.  The sector's structure
must be stated precisely, because a single adjective otherwise does
invisible work.

\begin{definition}[The gravitational source, exactly]
\label{def:gravsource}
The transverse-traceless spin-2 sector couples \emph{linearly} to the
matter stress words.  The longitudinal sector is not a degree of
freedom: it is fixed by its source through the gravitational Gauss law,
exactly as the Coulomb field is fixed by the charge, and the source of
the constraint is the total energy density \emph{including the field's
own}, carried exactly once by the coordinate-measure mass function
$m(r) = \int_0^r 4\pi r'^2 \rho\, dr'$: the field self-energy enters
through the \emph{measure} (coordinate volume in place of proper
volume), their difference being the gravitational binding energy
exactly.  No iterated dynamical self-stress vertex is declared, and no
nonlinear field equation is imported.  The clamped metric read is
$g_{00} = 1 - 2Gm(r)/r$ with $g_{rr} = 1/g_{00}$; the linearity is in
the variable $m$, not a linearization of anything.
\end{definition}

\begin{remark}[Relation to the self-coupling bootstrap]\label{rem:deser}
A reader trained on the Fierz--Pauli-to-Einstein bootstrap
\cite{feynmangrav,deser1970} will ask whether
Definition~\ref{def:gravsource} is that construction.  It is not, and
the difference matters in both directions.  The bootstrap iterates a
coupling of the linear field to a total stress that includes the
field's own; its conclusion is frequently cited as an unconditional
identity between linear spin-2 and general relativity, but the
iteration requires a choice of gravitational stress tensor at every
step (superpotential and field-redefinition freedom), and with that
freedom the bootstrap is a construction with inputs, not an identity
\cite{padmanabhan,deser2010}.  The present framework does not perform
the iteration and does not inherit the ambiguity: the field's
self-energy enters the source exactly once, through the solved
constraint of Definition~\ref{def:gravsource}, fixed by the same
quotient that fixes the Coulomb interaction: there is no stress
tensor to choose.  The comparison with general relativity is therefore
made emission by emission, and the first place the two constructions
could differ on a purely two-body observable is stated below as an open
problem rather than absorbed into a bootstrap citation.
\end{remark}

The ledger of the comparison, at the orders measured.  \emph{Test-mass
sector:} matter geodesics in the clamped read metric carry the
perihelion advance, light deflection, gravitational redshift, and
Shapiro delay at their measured orders.  Frame dragging is not carried
by that metric, which is static and diagonal and has no $g_{0i}$; it is
carried by the linear mass-current ($T^{0i}$) coupling of the same
spin-2 sector, the channel that also supplies the retardation term
below, and linear theory reproduces the Lense--Thirring precession at
measured accuracy.  \emph{Two-body
first post-Newtonian order:} the velocity structure of the
Einstein--Infeld--Hoffmann dynamics \cite{eih} is derived within the
sector as declared: the instantaneous exchange read gives
$(Gm_1m_2/2r)\,[\,3(v_1^2{+}v_2^2) - 8\,v_1\!\cdot\!v_2\,]$, and the
retardation kernel supplies exactly
$(Gm_1m_2/2r)\,[\,v_1\!\cdot\!v_2 - (n\!\cdot\!v_1)(n\!\cdot\!v_2)\,]$,
summing to the EIH form
$(Gm_1m_2/2r)\,[\,3(v_1^2{+}v_2^2) - 7\,v_1\!\cdot\!v_2 -
(n\!\cdot\!v_1)(n\!\cdot\!v_2)\,]$.  The computation is half a page and
is given in Appendix~\ref{app:eih}, together with its abelian control:
the identical kernel identity evaluated in the $U(1)$ sector yields the
Darwin Lagrangian \cite{landau}, a known result.  \emph{Two-body second
post-Newtonian order:} open, and stated as such.

\begin{openproblem}[The first purely two-body self-source observable]
\label{op:2pn}
At second post-Newtonian order the two-body dynamics first carries a
coefficient, linear in the symmetric mass ratio $\nu$, that no
test-mass read and no constraint normalization determines: the $-7\nu$ of the known two-body dynamics \cite{damourschafer}.  Derive it
within the sector as declared, from the forced-source structure of
Definition~\ref{def:gravsource}.  Until this derivation exists, the
gravitational sector's agreement with general relativity is
\emph{derived} through first post-Newtonian order and the test-mass
series, and \emph{required but not derived} at second post-Newtonian
two-body order; a derivation producing any value other than the known
one refutes the sector as declared.  The standard expectation is that a
linear field with a solved constraint fails at precisely this
coefficient; the framework does not dispute where the burden lies and
accepts the exposure as stated.  The problem's scope includes the mass
function of Definition~\ref{def:gravsource}: that the solved
constraint carries the binding energy through the measure is derived at
the leading order the comparisons above use, and asserted beyond it.
\end{openproblem}

Radiation reads are unchanged by the above: the binary-pulsar orbital
decay is the certified emission bracket of the spin-2 sector; a
coalescence is read from its radiation alone (radiated energy and
angular momentum are certified quadratic brackets on the emitted
transverse-traceless field, the final mass and spin are conservation
reads, the ringdown is the linear spectral read of the final field).
Because the interaction is attractive, the free-mass source problem is
unbounded below and the instrument fails shut on it; gravitational reads
enter through the clamped specialization, as they must for a
self-attracting source.

\subsection{An executed control}\label{sec:hydrogen}
The certified nonrelativistic hydrogen frequency ratio
$f(1S{-}3S)/f(1S{-}2S)$, emitted with certified width
$2.9 \times 10^{-12}$, excludes the measured value by $+1.309 \times
10^{-6}$; the exclusion is quantitatively the uncarried relativistic
sector: the Dirac correction $+1.109 \times 10^{-6}$ plus the
Lamb-shift differential.  What this exercise establishes, exactly: the
emission pipeline (interval arithmetic, the relational
calibration-free form of the ratio, and the exclusion semantics by
which the specialization brackets of the ladder are read) operates end
to end on a case where both numbers are independently known.  What it
does not establish: the ratio at this content is Rydberg-exact, so no
moment-cone certificate is exercised, and no claim of the theory is
tested that was not already known analytically.  In the paper's own
vocabulary this is a control, not a validation; the nontrivial
certificate machinery is exercised where the instance is not exactly
solvable, as in \S\ref{sec:instance}.

\subsection{The ultraviolet axis}\label{sec:uvaxis}
Every theorem of this section holds at stated budgets.  Whether
certificates exist whose constants, with the calibration anchors
re-solved at each ultraviolet budget, are uniform in the budget is the
refinement problem of \S\ref{sec:continuum} evaluated at this content:
renormalized limits of exactly this kind are proved for the Nelson model
\cite{nelson}, cutoff-uniform stability constants are proved for
nonrelativistic quantized-field matter \cite{ffg}, and the scalar
triviality results \cite{adc} neither establish nor preclude anything
here.  Along the axis the anchored relational emissions are certified
intervals, so convergence or drift of the renormalized family is
measured data; cutoff insensitivity is, on this instrument, a number
with a proof, or it is false.  The abelian ultraviolet axis and the
nonabelian chain of Conjecture~\ref{conj:chain} are one problem, the existence of certified contraction along refinement, posed at two
declared contents.

\section{Concluding remarks}\label{sec:conclusion}

The theory is a constructor: from declared content it builds the state,
and from the state its time, geometry, order, thermodynamics, and,
conditionally, its continuum, each object canonical, each emission a
verified two-sided bracket, each unproven step a named statement.  The
proved core comprises the kinematics and soundness results
(\S\ref{sec:kinematics}), the canonical constructions with the
factorization and harmonicity theorems (\S\ref{sec:canonical}), the
density theorem (\S\ref{sec:extensivity}), the exactness of refinement
and of two dimensions, the dichotomy, localization, no-go, and reduction
theorems (\S\ref{sec:continuum}), and the electrodynamic-matter
theorems of \S\ref{sec:subsumption}.  Proof statuses are graded in
place: complete proofs (the densities, tower summation, two dimensions,
uniformity, persistence, dichotomy, localization, no-decision-by-data,
clamped exactness, and the derivations of the appendices); complete
proofs under stated domain hypotheses (harmonicity, factorization);
and outlines with their deferred steps named (the finite-range
embedding; the closedness import in the duality lemma).  The abstract
claims completeness only where it holds.  The open surface comprises:
the fermionic machinery of
Open Problem~\ref{op:car}; the carrier transport of Open
Problem~\ref{op:carrier} and the explicit stability constants of
\S\ref{sec:gaps}; the local factorization for fermionic content of Open
Problem~\ref{op:sign}, which unlike the others may be a structural
obstruction rather than an unfinished construction; the dual-pair
presentation of Open Problem~\ref{op:dualpair}, with
Conjecture~\ref{conj:signfree} its candidate criterion for that
boundary; the convergence
completions of Open
Problem~\ref{op:qedconv}; the second post-Newtonian two-body derivation
of Open Problem~\ref{op:2pn}; Hypothesis~\ref{hyp:gma} with its four
decisive
free-fermion tests and the signature-consistency requirement they also
carry (Definition~\ref{def:signature}); the free-anchor slack computation of
\S\ref{sec:continuum}; and Conjecture~\ref{conj:chain}, whose statement
Theorems~\ref{thm:totaldichotomy} and \ref{thm:nogo} fence exactly: both
branches carry witnesses, no accumulation of certified data decides
between them, and resolution comes only by exhibition of a finite
verified object or by a structural monotone.  One decisive computation
the framework imposed on itself has been executed and passed: the Weyl
consistency comparison of Definition~\ref{def:weyl} on the worked
instance closes exactly, both sides in closed form
(Theorem~\ref{thm:weylinstance}; Appendix~\ref{app:weyl}), the first
decisive test of the geometric identification, passed at the strength
of an algebraic identity.  Four remain
unexecuted, and are
named as such: the prepotential re-pin of \S\ref{sec:correspondence},
the certified brackets of the worked instance re-derived in dual-pair
words with digit-level agreement required; the
free-anchor slack computation of \S\ref{sec:continuum}; the four
free-fermion tests of Hypothesis~\ref{hyp:gma}, blocked on Open
Problem~\ref{op:car}; and the signature-consistency comparison of
Definition~\ref{def:signature}, downstream of the same tests, on which
the real form of the spacetime group is an output.  A further grading concerns provenance: several
labeled theorems are standard facts transported into this frame and
cited as such (the density theorem is a subadditivity argument; the
canonical stack recombines the GNS, Tomita--Takesaki, and
Cipriani--Sauvageot constructions; the fusion and closure propositions
of \S\ref{sec:correspondence} are classical invariant theory and Howe
duality); they are numbered for the
dependency graph, not claimed as new.  What the paper offers as its own
are the resolving-emission and exclusion discipline, the factorization
scope with the sign problem in carrier form, the exactness of
refinement, the total dichotomy with its no-decision corollary, the
reduction theorem, the field--gauge correspondence read as
group-from-closure with its dual-pair criterion and its spacetime
reading (complex point from the ideal, real form from the modular
conjugations), and the quotient-fixed
gravitational reads with
their stated exposure.  The calibration requirements of
\S\ref{sec:continuum} likewise bind; the hadronic targets of
\S\ref{sec:hadrons} assemble from the theorems and open statements
above with nothing additional, which is the intended sense in which
this document presents one theory rather than a collection of
results.

\appendix

\section{Derivations for the worked instance}\label{app:derivations}
Charts appear below as evaluation scaffolding for chart-free quantities
(a Hausdorff volume, a distance, a spectral expansion); they enter no
definition and no emission.

\subsection{The metric}\label{app:geometry}
On one copy of SU(2) write $U = \cos\theta + i\sin\theta\,\hat n \cdot
\sigma$, $\theta \in [0,\pi]$, $\hat n \in S^2$.  Bi-invariance forces
the round form $ds^2 = R^2[\,d\theta^2 + \sin^2\!\theta\,
|d\hat n|^2\,]$, and the normalization follows from the verified
carr\'e du champ: with $x = 2\cos\theta$,
$\Gamma(x,x) = g^{\theta\theta}(\partial_\theta x)^2 =
g^{\theta\theta}\, 4\sin^2\!\theta$ must equal $(4 - x^2)/4 =
\sin^2\!\theta$, so $g^{\theta\theta} = 1/4$, i.e.\ $g_{\theta\theta} =
R^2 = 4$.  On the invariant sector of the pair, use coordinates
$(\theta_1, \theta_2, \beta)$ with $\beta$ the angle between the two
axes.  Then $|\nabla \theta_i|^2 = 1/4$; the axis sphere of copy $i$
has radius $2\sin\theta_i$, so moving axis $i$ changes $\beta$ at rate
$|\nabla^{(i)}\beta|^2 = 1/(4\sin^2\!\theta_i)$; and the $\theta$- and
$\beta$-gradients are mutually orthogonal (distinct tensor factors;
radial versus tangential on each factor).  Hence
\[
g^{\beta\beta} = \frac{s_1^2 + s_2^2}{4\, s_1^2 s_2^2},
\qquad
g = \mathrm{diag}\!\Big(4,\ 4,\ \frac{4\, s_1^2 s_2^2}{s_1^2 +
s_2^2}\Big),
\qquad
\sqrt{\det g} = \frac{8\, s_1 s_2}{\sqrt{s_1^2 + s_2^2}},
\]
with $s_i = \sin\theta_i$.

\subsection{The volume}\label{app:volume}
\[
\mathrm{Vol} = \int_0^\pi\! d\beta \int_0^\pi\!\!\int_0^\pi
\frac{8\, s_1 s_2}{\sqrt{s_1^2 + s_2^2}}\, d\theta_1\, d\theta_2
= 8\pi \int_{[-1,1]^2} \frac{du\, dv}{\sqrt{2 - u^2 - v^2}}
\]
(substituting $u = \cos\theta_1$, $v = \cos\theta_2$).  Split the
square into the inscribed unit disk and four corner regions.  Disk:
$\int_0^{2\pi}\!\!\int_0^1 \rho\, d\rho\, d\varphi\,(2 -
\rho^2)^{-1/2} = 2\pi(\sqrt2 - 1)$.  Corner (first quadrant, $u^2 + v^2
\ge 1$): in polar form
$J = 2\int_0^{\pi/4}\big[1 - \sqrt{\cos 2\varphi}/\cos\varphi\big]
d\varphi$; the substitutions $s = \sin\varphi$, then $s =
\sin\psi/\sqrt2$, give $\int_0^{\pi/4}\sqrt{\cos2\varphi}/\cos\varphi\,
d\varphi = \sqrt2\int_0^{\pi/2}\!\cos^2\!\psi/(1+\cos^2\!\psi)\,d\psi =
(\pi/2)(\sqrt2 - 1)$, hence $J = \pi(3/2 - \sqrt2)$.  The square
integral is $2\pi(\sqrt2 - 1) + 4\pi(3/2 - \sqrt2) = (4 - 2\sqrt2)\pi$,
and
\[
\mathrm{Vol} = 8\pi \cdot (4 - 2\sqrt2)\pi = (32 - 16\sqrt2)\,\pi^2 .
\]

\subsection{Distances}\label{app:distances}
One loop: class functions of a single copy carry the metric
$4\,d\theta^2$, so
\[
d = 2|\theta_1 - \theta_2|
  = 2|\arccos(x_1/2) - \arccos(x_2/2)|.
\]
Diameter of the two-loop variety: the projection
$(\theta_1,\theta_2,\beta) \mapsto (\theta_1,\theta_2)$ is
metric-nonincreasing onto the flat square with metric $4(d\theta_1^2 +
d\theta_2^2)$, giving
\[
d(\mathbf 1, (-\mathbf 1,-\mathbf 1))
  \ge 2\sqrt{\pi^2 + \pi^2} = 2\sqrt2\,\pi;
\]
the constant-$\beta$ diagonal
path attains this length, so equality holds.  That this pair realizes
the diameter follows from the same comparison run upward: the
invariant sector is a metric quotient of $S^3(2)\times S^3(2)$
(radius-$2$ round factors), quotient maps are distance-nonincreasing,
and the product's diameter is
\[
\sqrt{(2\pi)^2 + (2\pi)^2} = 2\sqrt2\,\pi,
\]
so no pair is farther.

\subsection{The spectral side and Weyl consistency}\label{app:weyl}
This subsection proves Theorem~\ref{thm:weylinstance}.  The general
statement behind it is classical: the equivariant Weyl law of
Donnelly and Br\"uning--Heintze \cite{donnelly,bruningheintze},
by which the counting function of invariant eigenfunctions is
governed by the volume and dimension of the orbit space.  What is
verified here is therefore not the law but the framework: every
normalization on both sides was fixed elsewhere, and the instance is
computed directly in closed form.  The carrier at
the reference is the closure of the invariant algebra in
$L^2(\mathrm{SU}(2)^2, \mathrm{Haar})$: the
simultaneous-conjugation-invariant sector.

\emph{Multiplicities.}  Peter--Weyl gives $L^2(\mathrm{SU}(2)) =
\bigoplus_j V_j \otimes V_j^*$, and conjugation acts on the $j$ block
as the adjoint-type action on $\mathrm{End}(V_j) \cong V_j \otimes
V_j^* \cong \bigoplus_{l=0}^{2j} V_l$ (integer $l$).  For the pair, the
diagonal-conjugation invariants of the $(j_1,j_2)$ block number
$\dim\mathrm{Inv}\big(\textstyle\bigoplus_{l_1} V_{l_1} \otimes
\bigoplus_{l_2} V_{l_2}\big) = \#\{l : l \le \min(2j_1, 2j_2)\} =
\min(2j_1,2j_2)+1$, at Casimir eigenvalue $j_1(j_1{+}1) +
j_2(j_2{+}1)$.  Cross-check on the low levels: the values $0, \tfrac34,
\tfrac32, 2$ carry multiplicities $1, 2, 2, 2$ (the words $\mathbf 1$;
$x, y$; $z, xy$; $\chi_1(U_1), \chi_1(U_2)$), matching the spectrum
used in the weak-coupling tower of \S\ref{sec:instance}.

\emph{Exact counting function.}  Substituting $a = 2j_1{+}1$,
$b = 2j_2{+}1 \in \mathbb N$, the eigenvalue is $(a^2 + b^2 - 2)/4$
with multiplicity $\min(a,b)$, so
\[
N(E) \;=\; \sum_{a,b \ge 1,\; a^2 + b^2 \le 4E + 2} \min(a,b),
\]
exact arithmetic with no numerical step.

\emph{Asymptotics.}  With $R^2 = 4E + 2$, comparison of the lattice sum
with its integral costs $O(R^2)$ (the summand is bounded by $R$ and
varies by at most $1$ per unit cell over a domain of $O(R^2)$ cells),
and in polar coordinates
\[
\iint_{a,b > 0,\ a^2+b^2 \le R^2} \min(a,b)\, da\, db
= \frac{R^3}{3}\int_0^{\pi/2} \min(\cos\varphi, \sin\varphi)\, d\varphi
= \frac{R^3}{3}\,(2 - \sqrt2),
\]
since $\int_0^{\pi/2}\min(\cos,\sin) = 2\int_0^{\pi/4}\sin = 2 -
\sqrt2$.  With $R^3 = 8E^{3/2}(1 + O(1/E))$,
\[
N(E) \;=\; \frac{16 - 8\sqrt2}{3}\, E^{3/2} \;+\; O(E).
\]
The limit exists, so Dixmier measurability holds for this instance and
the spectral volume is admissible as an emission
(\S\ref{sec:metricdim}).

\emph{The identity.}  A three-dimensional geometry of volume $V$ has
Weyl coefficient $V \omega_3/(2\pi)^3 = V/(6\pi^2)$, so the volume read
from the count is $6\pi^2 \cdot (16 - 8\sqrt2)/3 = 2\pi^2(16 -
8\sqrt2) = (32 - 16\sqrt2)\pi^2$: exactly the closed form of
Appendix~\ref{app:volume}.  The $\sqrt2$ arises on the metric side
from the corner integral of Appendix~\ref{app:volume} and on the
spectral side from $\int_0^{\pi/2}\min(\cos,\sin)\,d\varphi$: two
unrelated computations forced to one number.

\emph{Coupling independence.}  For $h = \mathrm{Cas} - \lambda M$ with
$\|M\| \le 6$, the operator inequalities $\mathrm{Cas} - 6\lambda \le
h \le \mathrm{Cas} + 6\lambda$ give $N_{\mathrm{Cas}}(E - c) \le
N_{h - E_0}(E) \le N_{\mathrm{Cas}}(E + c)$ with $c = 6\lambda +
|E_0|$; bounded shifts leave the $E^{3/2}$ coefficient unchanged, so
Theorem~\ref{thm:weylinstance} holds at every declared coupling.

\begin{remark}[One-loop warm-up; finite-$E$ diagnostic]
\label{rem:weylnumerics}
The single-plaquette instance closes the same way in dimension one:
eigenvalues $j(j{+}1)$ with multiplicity one give $N(E) = \#\{k \ge 1 :
k^2 \le 4E{+}1\} \sim 2\sqrt E$, and the one-dimensional Weyl law
$N \sim (L/\pi)\sqrt E$ reads off $L = 2\pi$, the length of $[0,\pi]$
in the metric $4\,d\theta^2$, matching the one-loop distance formula of
\S\ref{sec:geometry}.  As a diagnostic of the subleading (boundary)
term in the two-loop count, the exact sum at $E = 10^6$ gives
$N(E)/E^{3/2} = 1.5621140$ against the limit $(16-8\sqrt2)/3 =
1.5620972$, the residual carried by the $O(E)$ term.
\end{remark}

\subsection{Strong coupling}\label{app:strong}
Near the identity class use geodesic normal coordinates $s_i \in
\mathbb R^3$ per loop ($\theta_i = |s_i|/2$ from $g_{\theta\theta} =
4$).  Then $x \approx 2 - |s_1|^2/4$, $y \approx 2 - |s_2|^2/4$, and,
since small rotations compose by adding their vectors at leading order,
$z \approx 2 - |s_1 + s_2|^2/4$; hence
\[
6 - M \approx \tfrac14\big(|s_1|^2 + |s_2|^2 + |s_1 + s_2|^2\big),
\qquad
h + 6\lambda \approx -\Delta_{s_1} - \Delta_{s_2} + \tfrac{\lambda}{4}
\big(s_1^2 + s_2^2 + (s_1{+}s_2)^2\big),
\]
with the flat Laplacians of the $\Gamma$-metric.  Per component the
stiffness matrix is $\tfrac{\lambda}{4}
\big(\begin{smallmatrix}2&1\\1&2\end{smallmatrix}\big)$ with
eigenvalues $3\lambda/4$ and $\lambda/4$; normal modes $u = (s_1 +
s_2)/\sqrt2$, $v = (s_1 - s_2)/\sqrt2$, three components each; the
ground energy of $-d^2/ds^2 + k s^2$ is $\sqrt k$ (spectrum
$(2n{+}1)\sqrt k$: level spacing $2\sqrt k$, twice the zero-point
energy), so
\[
E_1 + 6\lambda = 3\sqrt{3\lambda/4} + 3\sqrt{\lambda/4} =
\tfrac32(1 + \sqrt3)\sqrt{\lambda} + O(1),
\]
the $O(1)$ collecting measure (drift) and quartic corrections,
computed exactly below.  Excitations are
priced by the level spacing $\omega = 2\sqrt k$, not the zero-point
energy $\sqrt k$.  Single quanta of $u$ or $v$ are vectors under the
residual diagonal SO(3) and do not lie in the invariant sector, whose
functions are built from $|s_1|^2$, $|s_2|^2$, $s_1 \cdot s_2$; the
lowest invariant excitation is therefore a soft-mode pair at
$2\omega_v = 4\sqrt{\lambda/4} = 2\sqrt\lambda$, and the next is the
mixed scalar $u \cdot v$ at $\omega_u + \omega_v =
(1 + \sqrt3)\sqrt\lambda$.  Hence the gap-ratio asymptote
$(1 + \sqrt3)/2$ (a ratio of spacings, insensitive to the pricing),
and, since $v$ is odd and $v^2$ even under loop
exchange, the swap-even ordering at strong coupling.  The virial
theorem gives $\lambda\langle 6 - M\rangle = E_{\mathrm{harm}}/2$,
i.e.\ $\langle M\rangle = 6 - \tfrac34(1 + \sqrt3)/\sqrt\lambda +
O(1/\lambda)$.

\paragraph{The constant term.}  The $O(1)$ term is fixed by an exact
conjugation.  Per SU(2) factor, with $J(s) = [\sin(|s|/2)/(|s|/2)]^2$
the Jacobian of the exponential chart, the Casimir satisfies
\[
C_2 \;=\; J^{-1/2}\,\big(-\Delta_{\mathrm{flat}} -
\tfrac14\big)\,J^{1/2}
\]
\emph{exactly}: $J^{1/2}\chi_j = 2\sin\!\big((2j{+}1)|s|/2\big)/|s|$ is
a flat radial eigenfunction at $((2j{+}1)/2)^2 = j(j{+}1) + \tfrac14$
for every $j$ (verified symbolically through $2j = 6$).  Hence
$h + 6\lambda$ is unitarily a \emph{flat} six-dimensional problem,
\[
h + 6\lambda \;\cong\; -\Delta_6 - \tfrac12 + \lambda W, \qquad
W = 6 - M \text{ in normal coordinates},
\]
with $W$ exact trigonometry: from $\operatorname{tr}(U_1U_2) =
2[\cos u_1 \cos u_2 - (\hat s_1 \!\cdot \hat s_2)\sin u_1 \sin u_2]$,
$u_i = |s_i|/2$,
\[
W = \tfrac14\big(|s_1|^2 + |s_2|^2 + |s_1{+}s_2|^2\big)
- \frac{|s_1|^4 + |s_2|^4}{96} - \frac{|s_1|^2|s_2|^2}{32}
- \frac{(s_1\!\cdot s_2)(|s_1|^2 + |s_2|^2)}{48} + O(s^6),
\]
the quadratic part reproducing the stiffnesses above.  In the normal
modes $u = (s_1{+}s_2)/\sqrt2$, $v = (s_1{-}s_2)/\sqrt2$ the harmonic
ground state is the product Gaussian with per-component variances
$\sigma_u^2 = 1/\sqrt{3\lambda}$, $\sigma_v^2 = 1/\sqrt{\lambda}$, and
first-order perturbation theory in the quartic gives
$\lambda\langle W_4\rangle = -(30 + 13\sqrt3)/192$; second order in
$W_4$ and first order in $W_6$ are both $O(\lambda^{-1/2})$.  With the
exact $-\tfrac12$ from the conjugation,
\[
c_0 \;=\; -\tfrac12 - \frac{30 + 13\sqrt3}{192}
\;=\; -\frac{126 + 13\sqrt3}{192} \;\approx\; -0.77352.
\]
Controls: (i) the one-loop analogue of the identical mechanics
predicts $c_0^{(1)} = -\tfrac14 - \tfrac{5}{64} = -\tfrac{21}{64}$,
and the exact character-basis tridiagonal reproduces it to $10^{-9}$
under a large-$\lambda$ fit; (ii) every algebraic step (trace
identity, quartic expansion, Gaussian moments, conjugation
eigenfunctions) verifies symbolically; (iii) converged exact-moment
variational values at $\lambda \in \{2,4,8\}$ sit within $0.014$ of
$c_0$ with subleading coefficients of natural size, and, with $c_0$
fixed, the implied $(c_{-1}, c_{-2})$ postdict the certified
$\lambda = 1/2$ window; (iv) the derived value lies inside the
certified $\lambda = 2$ window outright.

\subsection{The fourth-order weak-coupling coefficient}\label{app:fourth}
At $\lambda = 0$ the ground state is the Haar state with $E_0 = 0$ and
$\langle M\rangle = 0$.  The span of invariant monomials of degree
$\le 2$ (ten elements) is invariant under the Casimir, closing by the
$\Gamma$-calculus of \S\ref{sec:instance}, and decomposes exactly with
eigenvalues $0, \tfrac34, \tfrac32, 2, \tfrac{11}4, 4$ and
multiplicities $1,2,2,2,2,1$, matching the combinatorial audit
$\min(2j_1,2j_2)+1$ of Appendix~\ref{app:weyl}.  Every intermediate
state of fourth-order Rayleigh--Schr\"odinger perturbation theory lies
in this span (each application of $M$ raises monomial degree by at
most one and the path must return to the vacuum), so the fourth order
is a finite computation in exact rational arithmetic with the reduced
resolvent $R = -\sum_{\mu \ne 0} P_\mu/\mu$:
\[
E^{(4)} \;=\; \langle 0|VRVRVRV|0\rangle - E^{(2)}\,\langle
0|VR^2V|0\rangle
\;=\; -\tfrac{3409}{297} - \big(-\tfrac{10}{3}\big)(4)
\;=\; \tfrac{551}{297},
\]
with $V = -M$ per unit coupling and $\langle 0|VR^2V|0\rangle = 4$
exactly.  The same machinery reproduces $E^{(2)} = -\tfrac{10}{3}$ and
$E^{(3)} = -\tfrac{32}{9}$, and an independent variational route
(exact-moment Rayleigh--Ritz at monomial degree $\le 6$, couplings
$\lambda \in [1/96, 1/32]$) confirms the value.

The fifth order follows from the same objects by the $2n{+}1$ rule:
with $\psi_1 = RV|0\rangle$ and $\psi_2 = RVRV|0\rangle$ (both inside
the degree-$\le 2$ span, so no further resolvent is needed),
\[
E^{(5)} = \langle\psi_2|V|\psi_2\rangle
- 2E^{(2)}\langle\psi_1|\psi_2\rangle
- E^{(3)}\langle\psi_1|\psi_1\rangle
= \tfrac{172064}{9801} \;\approx\; +17.5558,
\]
confirmed variationally: with $c_4$ fixed, the residual fit gives
$c_5 = 17.579$, and the diagnostic ratio
$\mathrm{resid}/\lambda^5$ marches monotonically onto the exact value
as $\lambda \to 0$.

A caution on reading coefficients off bracket data: the higher-order
coefficients here are large, so naive fits mislead---a two-point fit
at $\lambda \approx 1/16$ omitting the fifth order reports $c_4 +
c_5\lambda \approx 3$ rather than $c_4 \approx 1.86$, and a joint fit
of $(c_4, c_5)$ leaks $c_6 \approx 21$ into the fifth-order estimate.
Certified brackets bound energies, not series coefficients; the
coefficients must be derived, with the brackets as their check.

\section{The first post-Newtonian two-body read}\label{app:eih}
The velocity structure derived below was first obtained from
one-graviton exchange by Iwasaki \cite{iwasaki}; the derivation is
retained because the sector as declared (forced source, no dynamical
self-stress vertex, Definition~\ref{def:gravsource}) must be shown
to produce it, and because the abelian control pins the conventions.
Between conserved sources the single-exchange kernel is unique at this
order (gauge-dependent parts of the propagator drop against
conservation); the propagator algebra below is evaluation scaffolding
for a chart-free read, the two-body interaction energy.

\emph{Instantaneous part.}  Worldline stress $T_i^{\mu\nu} = m_i\,
u_i^\mu u_i^\nu/\gamma_i$ (the $1/\gamma$ from the worldline density);
spin-2 numerator $P_{\mu\nu\rho\sigma} = \tfrac12(\eta_{\mu\rho}
\eta_{\nu\sigma} + \eta_{\mu\sigma}\eta_{\nu\rho} -
\eta_{\mu\nu}\eta_{\rho\sigma})$, so $T_1 P\, T_2 = \tfrac12 m_1 m_2
[\,2(u_1 \cdot u_2)^2 - 1\,]/(\gamma_1\gamma_2)$.  Normalizing the
static term to $Gm_1m_2/r$ and expanding $u_1 \cdot u_2 =
\gamma_1\gamma_2(1 - v_1 \cdot v_2)$:
\[
\begin{aligned}
\frac{2(u_1 \cdot u_2)^2 - 1}{\gamma_1\gamma_2}
  &= 1 + \tfrac32(v_1^2 + v_2^2)
     - 4\, v_1 \cdot v_2 + O(v^4), \\
L_{\mathrm{inst}}
  &= \frac{Gm_1m_2}{r}
     + \frac{Gm_1m_2}{2r}
       \big[\,3(v_1^2 + v_2^2) - 8\, v_1 \cdot v_2\,\big].
\end{aligned}
\]

\emph{Retardation.}  $1/(k^2 - \omega^2) = 1/k^2 + \omega^2/k^4 +
O(\omega^4)$, with $1/k^2 \leftrightarrow 1/(4\pi r)$ and $1/k^4
\leftrightarrow -\,r/(8\pi)$.  The $\omega^2$ term between the static
channels is $\Delta L = -\tfrac{G m_1 m_2}{2}\,\partial_t
\partial_{t'}\, r(t,t')\big|_{t'=t}$ with $r(t,t') = |x_1(t) -
x_2(t')|$; since $\partial_t \partial_{t'} r = -\big[v_1 \cdot v_2 -
(n \cdot v_1)(n \cdot v_2)\big]/r$,
\[
\Delta L = +\frac{Gm_1m_2}{2r}\big[\,v_1 \cdot v_2 -
(n \cdot v_1)(n \cdot v_2)\,\big].
\]

\emph{Sum.}
\[
L_{\mathrm{inst}} + \Delta L = \frac{Gm_1m_2}{r} +
\frac{Gm_1m_2}{2r}\big[\,3(v_1^2 + v_2^2) - 7\, v_1 \cdot v_2 -
(n \cdot v_1)(n \cdot v_2)\,\big],
\]
the Einstein--Infeld--Hoffmann velocity structure \cite{eih}.

\emph{Abelian control.}  The same two lines in the $U(1)$ sector:
current numerator $(u_1 \cdot u_2)/(\gamma_1\gamma_2) = 1 - v_1 \cdot
v_2$ exactly, static sign reversed, retardation kernel identical.  Then
$L = -q_1q_2/r + (q_1q_2/r)\, v_1 \cdot v_2 - (q_1q_2/2r)[\,v_1 \cdot
v_2 - (n \cdot v_1)(n \cdot v_2)\,] = -q_1q_2/r +
(q_1q_2/2r)[\,v_1 \cdot v_2 + (n \cdot v_1)(n \cdot v_2)\,]$: the
Darwin Lagrangian \cite{landau}, a known result, validating the kernel
identity and the sign conventions.

\section{Executed falsifier battery}\label{app:battery}
Every check computes a ground truth by an independent route (exact
diagonalization, matrix logarithm, or Fock-space projection) and compares
it to the closed form; no empirical input enters.  The closures instantiate
Theorems~\ref{thm:gapid}--\ref{thm:closure}: where a bracket closes to
$L^* = U^*$ the certificate factors annihilate the state, and that
annihilation is verified.

\textbf{Free-field zero.}  The level-1 optimizer of $\Tr(h\gamma)$ is the
negative-energy projector (the Dirac sea as output) to $0$; the subtracted
vacuum energy is $0$; the charge-graded antiparticle energies are strictly
positive (vacuum stable), catching any grading-sign error.

\textbf{van Hove.}  The completed square of
Definition~\ref{def:bosoncalc} is an operator identity to
$3.6\times10^{-14}$, positive-semidefinite, and the exact-diagonalization
ground energy meets the certificate value $-\sum|B|^2/\omega$ to $0$: the
boson word calculus verified, and tight ($L^*=U^*$).

\textbf{Modular reads.}  For a two-mode Gaussian state built by a Gaussian
unitary on a thermal state (no truncation assumption), the modular
Hamiltonian extracted from $-\log\rho$ is quadratic (evaluability), its
symplectic-log relation $\nu=\tfrac12\coth(g/2)$ holds, and the state is
KMS at inverse modular temperature one; all three converge with Fock
cutoff ($2\times10^{-3}\to6\times10^{-5}$), confirming physics rather than
artifact.  The rank-deficiency fail-shut and the two-mode Gaussian modular
flow (KMS to machine precision) are the temporal end-to-end pins.

\textbf{Two-mode Gaussian conditional (exact, no truncation).}  (a) The
total-charge conditional expectation commutes with the per-covolume
average to $0.0$; a symmetry-breaking per-mode conditioning fails to
commute at $0.41$, so the commutation test discriminates within the
conditioning class admitted by Definition~\ref{def:pip}.  (b) The
charge-conditioned two-point function
from the natural-orbital closed form equals the Fock ground truth to
$1.4\times10^{-16}$, and the certified two-sided bracket closes on the
conditional ground energy, $L^*=U^*=e_0$, certificate gap $0$: boundary
contact (Theorem~\ref{thm:closure}) exhibited.

\textbf{Emergence (harmonic half-pin).}  Zero-point energy exact by two
routes; the structure-factor read develops Bragg atoms at the
reciprocal-lattice vectors with the correct Debye--Waller intensity ratio.
Half-pin: the lattice is in the coefficients, so this tests the read, not
order-from-disorder, which needs the interacting engine.

\textbf{Cost-Theorem solvable models.}  RPA per-mode saturation (the
variational bound saturates the exact per-mode gap; totals
$\sim r_0^{-3}$); two-mode order structure (first-order vanishing);
three-mode commutator fluctuation ($O(\xi^2 s^2)$); four-mode conditioning
law ($A=2\lambda^2/\omega_2$, flat reduced densities have zero f-sum so
conditioning is load bearing); screened incoherent weight
($\sim r_0^{-6}$); the harmonic crystal and non-adiabatic H$_2$/HD massive
pins; van Hove, the harmonic charge, and spin--boson/Rabi radiative pins.

\begin{thebibliography}{99}
\bibitem{pusz} W.~Pusz, S.~L.~Woronowicz, \emph{Passive states and KMS
states for general quantum systems}, Comm. Math. Phys. \textbf{58}
(1978) 273--290.
\bibitem{coleman} A.~J.~Coleman, \emph{Structure of fermion density
matrices}, Rev. Mod. Phys. \textbf{35} (1963) 668--686.
\bibitem{mazziotti} D.~A.~Mazziotti, \emph{Realization of quantum
chemistry without wave functions through first-order semidefinite
programming}, Phys. Rev. Lett. \textbf{93} (2004) 213001.
\bibitem{npa} M.~Navascu\'es, S.~Pironio, A.~Ac\'in, \emph{A convergent
hierarchy of semidefinite programs characterizing the set of quantum
correlations}, New J. Phys. \textbf{10} (2008) 073013.
\bibitem{lin} H.~W.~Lin, \emph{Bootstraps to strings: solving random
matrix models with positivity}, J. High Energy Phys. \textbf{06} (2020)
090.
\bibitem{hhk} X.~Han, S.~A.~Hartnoll, J.~Kruthoff, \emph{Bootstrapping
matrix quantum mechanics}, Phys. Rev. Lett. \textbf{125} (2020) 041601.
\bibitem{andersonkruczenski} P.~D.~Anderson, M.~Kruczenski, \emph{Loop
equations and bootstrap methods in the lattice}, Nucl. Phys. B
\textbf{921} (2017) 702--726.
\bibitem{kazakovzheng} V.~Kazakov, Z.~Zheng, \emph{Bootstrap for lattice
Yang--Mills theory}, Phys. Rev. D \textbf{107} (2023) L051501.
\bibitem{kazakovzhengfinite} V.~Kazakov, Z.~Zheng, \emph{Bootstrap for
finite $N$ lattice Yang--Mills theory}, J. High Energy Phys.
\textbf{03} (2025) 099.
\bibitem{guosu3} Y.~Guo, Z.~Li, G.~Yang, G.~Zhu, \emph{Bootstrapping
SU(3) lattice Yang--Mills theory}, J. High Energy Phys. \textbf{12}
(2025) 033.
\bibitem{lizhou} Z.~Li, S.~Zhou, \emph{Bootstrapping the abelian
lattice gauge theories}, J. High Energy Phys. \textbf{08} (2024) 154.
\bibitem{lindzero} H.~W.~Lin, \emph{Bootstrap bounds on D0-brane
quantum mechanics}, J. High Energy Phys. \textbf{06} (2023) 038.
\bibitem{linzhengmatrix} H.~W.~Lin, Z.~Zheng, \emph{Bootstrapping
ground state correlators in matrix theory.  Part I}, J. High Energy
Phys. \textbf{01} (2025) 190.
\bibitem{kull} I.~Kull, N.~Schuch, B.~Dive, M.~Navascu\'es,
\emph{Lower bounds on ground-state energies of local Hamiltonians
through the renormalization group}, Phys. Rev. X \textbf{14} (2024)
021008.
\bibitem{fawzi} H.~Fawzi, O.~Fawzi, S.~O.~Scalet, \emph{Certified
algorithms for equilibrium states of local quantum Hamiltonians}, Nat.
Commun. \textbf{15} (2024) 7394.
\bibitem{peyrlparrilo} H.~Peyrl, P.~A.~Parrilo, \emph{Computing sum of
squares decompositions with rational coefficients}, Theoret. Comput.
Sci. \textbf{409} (2008) 269--281.
\bibitem{donnelly} H.~Donnelly, \emph{G-spaces, the asymptotic splitting
of $L^2(M)$ into irreducibles}, Math. Ann. \textbf{237} (1978) 23--40.
\bibitem{bruningheintze} J.~Br\"uning, E.~Heintze,
\emph{Representations of compact Lie groups and elliptic operators},
Invent. Math. \textbf{50} (1978) 169--203.
\bibitem{iwasaki} Y.~Iwasaki, \emph{Quantum theory of gravitation vs.\
classical theory: fourth-order potential}, Prog. Theor. Phys.
\textbf{46} (1971) 1587--1609.
\bibitem{migdal} A.~A.~Migdal, \emph{Recursion equations in gauge field
theories}, Sov. Phys. JETP \textbf{42} (1975) 413.
\bibitem{fot} M.~Fukushima, Y.~Oshima, M.~Takeda, \emph{Dirichlet Forms
and Symmetric Markov Processes}, 2nd ed., de Gruyter (2011).
\bibitem{guoli} Y.~Guo, W.~Li, \emph{Solving anharmonic oscillator with
null states: Hamiltonian bootstrap and Dyson--Schwinger equations},
Phys. Rev. D \textbf{108} (2023) 125002.
\bibitem{bdot} S.~Bravyi, D.~P.~DiVincenzo, R.~I.~Oliveira,
B.~M.~Terhal, \emph{The complexity of stoquastic local Hamiltonian
problems}, Quantum Inf. Comput. \textbf{8} (2008) 361--385.
\bibitem{mlh} M.~Marvian, D.~A.~Lidar, I.~Hen, \emph{On the
computational complexity of curing non-stoquastic Hamiltonians}, Nat.
Commun. \textbf{10} (2019) 1571.
\bibitem{howe} R.~Howe, \emph{Remarks on classical invariant theory},
Trans. Amer. Math. Soc. \textbf{313} (1989) 539--570.
\bibitem{weylbook} H.~Weyl, \emph{The Classical Groups: Their Invariants
and Representations}, Princeton University Press (1939).
\bibitem{collinssniady} B.~Collins, P.~\'Sniady, \emph{Integration with
respect to the Haar measure on unitary, orthogonal and symplectic
group}, Comm. Math. Phys. \textbf{264} (2006) 773--795.
\bibitem{schwinger} J.~Schwinger, \emph{On angular momentum}, in
\emph{Quantum Theory of Angular Momentum}, eds. L.~C.~Biedenharn,
H.~van~Dam, Academic Press (1965) 229--279.
\bibitem{mathur} M.~Mathur, \emph{Harmonic oscillator pre-potentials in
SU(2) lattice gauge theory}, J. Phys. A \textbf{38} (2005)
10015--10026.
\bibitem{anishetty} R.~Anishetty, M.~Mathur, I.~Raychowdhury,
\emph{Prepotential formulation of SU(3) lattice gauge theory}, J. Phys.
A \textbf{43} (2010) 035403.
\bibitem{quantumlink} R.~Brower, S.~Chandrasekharan, U.-J.~Wiese,
\emph{QCD as a quantum link model}, Phys. Rev. D \textbf{60} (1999)
094502.
\bibitem{jones} V.~F.~R.~Jones, \emph{Index for subfactors}, Invent.
Math. \textbf{72} (1983) 1--25.
\bibitem{deligne} P.~Deligne, \emph{La cat\'egorie des repr\'esentations
du groupe sym\'etrique $S_t$, lorsque $t$ n'est pas un entier naturel},
in \emph{Algebraic Groups and Homogeneous Spaces}, TIFR Mumbai (2007)
209--273.
\bibitem{voiculescu} D.~Voiculescu, \emph{Limit laws for random matrices
and free products}, Invent. Math. \textbf{104} (1991) 201--220.
\bibitem{takesaki} M.~Takesaki, \emph{Theory of Operator Algebras II},
Springer (2003).
\bibitem{cipriani} F.~Cipriani, J.-L.~Sauvageot, \emph{Derivations as
square roots of Dirichlet forms}, J. Funct. Anal. \textbf{201} (2003)
78--120.
\bibitem{connes} A.~Connes, \emph{Noncommutative Geometry}, Academic
Press (1994).
\bibitem{connesrecon} A.~Connes, \emph{On the spectral characterization
of manifolds}, J. Noncommut. Geom. \textbf{7} (2013) 1--82.
\bibitem{wiesbrock} H.-W.~Wiesbrock, \emph{Half-sided modular inclusions
of von Neumann algebras}, Comm. Math. Phys. \textbf{157} (1993) 83--92.
\bibitem{summers} D.~Buchholz, O.~Dreyer, M.~Florig, S.~J.~Summers,
\emph{Geometric modular action and spacetime symmetry groups}, Rev.
Math. Phys. \textbf{12} (2000) 475--560.
\bibitem{guidolongo} D.~Guido, R.~Longo, H.-W.~Wiesbrock, \emph{Extensions
of conformal nets and superselection structures}, Comm. Math. Phys.
\textbf{192} (1998) 217--244.
\bibitem{bisognano} J.~Bisognano, E.~Wichmann, \emph{On the duality
condition for a Hermitian scalar field}, J. Math. Phys. \textbf{16}
(1975) 985--1007.
\bibitem{knabe} S.~Knabe, \emph{Energy gaps and elementary excitations
for certain VBS-quantum antiferromagnets}, J. Stat. Phys. \textbf{52}
(1988) 627--638.
\bibitem{nachtergaele} B.~Nachtergaele, \emph{The spectral gap for some
spin chains with discrete symmetry breaking}, Comm. Math. Phys.
\textbf{175} (1996) 565--606.
\bibitem{bhm} S.~Bravyi, M.~Hastings, S.~Michalakis, \emph{Topological
quantum order: stability under local perturbations}, J. Math. Phys.
\textbf{51} (2010) 093512.
\bibitem{luscher} M.~L\"uscher, \emph{Volume dependence of the energy
spectrum in massive quantum field theories I}, Comm. Math. Phys.
\textbf{104} (1986) 177--206.
\bibitem{lanford} O.~E.~Lanford III, \emph{A computer-assisted proof of
the Feigenbaum conjectures}, Bull. Amer. Math. Soc. \textbf{6} (1982)
427--434.
\bibitem{koch} H.~Koch, P.~Wittwer, \emph{A non-Gaussian
renormalization-group fixed point for hierarchical scalar lattice field
theories}, Comm. Math. Phys. \textbf{106} (1986) 495--532.
\bibitem{uhlmann} A.~Uhlmann, \emph{Relative entropy and the
Wigner--Yanase--Dyson--Lieb concavity in an interpolation theory},
Comm. Math. Phys. \textbf{54} (1977) 21--32.
\bibitem{kosaki} H.~Kosaki, \emph{Relative entropy of states: a
variational expression}, J. Operator Theory \textbf{16} (1986) 335--348.
\bibitem{putinar} M.~Putinar, \emph{Positive polynomials on compact
semi-algebraic sets}, Indiana Univ. Math. J. \textbf{42} (1993)
969--984.
\bibitem{lasserre} J.-B.~Lasserre, \emph{Global optimization with
polynomials and the problem of moments}, SIAM J. Optim. \textbf{11}
(2001) 796--817.
\bibitem{niesch} J.~Nie, M.~Schweighofer, \emph{On the complexity of
Putinar's Positivstellensatz}, J. Complexity \textbf{23} (2007)
135--150.
\bibitem{powers} V.~Powers, C.~Scheiderer, \emph{The moment problem for
non-compact semialgebraic sets}, Adv. Geom. \textbf{1} (2001) 71--88.
\bibitem{marshall} M.~Marshall, \emph{Positive Polynomials and Sums of
Squares}, AMS Surveys \textbf{146} (2008).
\bibitem{bakry} D.~Bakry, M.~\'Emery, \emph{Diffusions hypercontractives},
S\'eminaire de probabilit\'es XIX, Springer LNM \textbf{1123} (1985)
177--206.
\bibitem{casini} H.~Casini, M.~Huerta, \emph{On the RG running of the
entanglement entropy of a circle}, Phys. Rev. D \textbf{85} (2012)
125016.
\bibitem{nielsen} H.~B.~Nielsen, M.~Ninomiya, \emph{Absence of neutrinos
on a lattice}, Nucl. Phys. B \textbf{185} (1981) 20--40.
\bibitem{ginsparg} P.~Ginsparg, K.~Wilson, \emph{A remnant of chiral
symmetry on the lattice}, Phys. Rev. D \textbf{25} (1982) 2649--2657.
\bibitem{luscherchiral} M.~L\"uscher, \emph{Exact chiral symmetry,
topological charge and related topics in lattice QCD}, Nucl. Phys. B
Proc. Suppl. \textbf{83} (2000) 34--43.
\bibitem{feynmangrav} R.~P.~Feynman, F.~B.~Morinigo, W.~G.~Wagner,
\emph{Feynman Lectures on Gravitation}, Addison--Wesley (1995).
\bibitem{deser1970} S.~Deser, \emph{Self-interaction and gauge
invariance}, Gen. Relativ. Gravit. \textbf{1} (1970) 9--18.
\bibitem{padmanabhan} T.~Padmanabhan, \emph{From gravitons to gravity:
myths and reality}, Int. J. Mod. Phys. D \textbf{17} (2008) 367--398.
\bibitem{deser2010} S.~Deser, \emph{Gravity from self-interaction
redux}, Gen. Relativ. Gravit. \textbf{42} (2010) 641--646.
\bibitem{eih} A.~Einstein, L.~Infeld, B.~Hoffmann, \emph{The
gravitational equations and the problem of motion}, Ann. Math.
\textbf{39} (1938) 65--100.
\bibitem{damourschafer} T.~Damour, G.~Sch\"afer, \emph{Higher-order
relativistic periastron advances and binary pulsars}, Nuovo Cimento B
\textbf{101} (1988) 127--176.
\bibitem{landau} L.~D.~Landau, E.~M.~Lifshitz, \emph{The Classical
Theory of Fields}, 4th ed., Pergamon (1975).
\bibitem{brownravenhall} G.~E.~Brown, D.~G.~Ravenhall, \emph{On the
interaction of two electrons}, Proc. R. Soc. Lond. A \textbf{208} (1951)
552--559.
\bibitem{lss} E.~H.~Lieb, H.~Siedentop, J.~P.~Solovej, \emph{Stability
and instability of relativistic electrons in classical electromagnetic
fields}, J. Stat. Phys. \textbf{89} (1997) 37--59.
\bibitem{liebloss} E.~H.~Lieb, M.~Loss, \emph{Stability of a model of
relativistic quantum electrodynamics}, Comm. Math. Phys. \textbf{228}
(2002) 561--588.
\bibitem{gll} M.~Griesemer, E.~H.~Lieb, M.~Loss, \emph{Ground states in
non-relativistic quantum electrodynamics}, Invent. Math. \textbf{145}
(2001) 557--595.
\bibitem{lieblebowitz} E.~H.~Lieb, J.~L.~Lebowitz, \emph{The
constitution of matter: existence of thermodynamics for systems composed
of electrons and nuclei}, Adv. Math. \textbf{9} (1972) 316--398.
\bibitem{fdll} C.~Fefferman, R.~de la Llave, \emph{Relativistic
stability of matter I}, Rev. Mat. Iberoam. \textbf{2} (1986) 119--213.
\bibitem{nelson} E.~Nelson, \emph{Interaction of nonrelativistic
particles with a quantized scalar field}, J. Math. Phys. \textbf{5}
(1964) 1190--1197.
\bibitem{ffg} C.~Fefferman, J.~Fr\"ohlich, G.~M.~Graf, \emph{Stability
of ultraviolet-cutoff quantum electrodynamics with non-relativistic
matter}, Comm. Math. Phys. \textbf{190} (1997) 309--330.
\bibitem{adc} M.~Aizenman, H.~Duminil-Copin, \emph{Marginal triviality
of the scaling limits of critical 4D Ising and $\varphi^4_4$ models},
Ann. of Math. \textbf{194} (2021) 163--235.
\end{thebibliography}

\end{document}
